\pdfoutput=1
\documentclass[11pt]{amsart}
\usepackage{amsmath,amssymb,amsthm}
\usepackage[margin=1.2in]{geometry}
\usepackage[colorlinks=true,linkcolor=blue,citecolor=blue,urlcolor=blue]{hyperref}
\hypersetup{pdfauthor={Bernd Johannes Wuebben},
 pdftitle={Extending the Kempf-Ness function to complexes of quiver representations},
 pdfsubject={MSC 2020: 16G20, 53D20; 14F08, 18G80, 53C07},
 pdfkeywords={quiver representation, complex of representations, stability condition, superconnection, polar decomposition, Kempf-Ness function, Donaldson functional}}
\setlength{\emergencystretch}{2em}
\newtheorem{theorem}{Theorem}[section]
\newtheorem{proposition}[theorem]{Proposition}
\newtheorem{lemma}[theorem]{Lemma}
\newtheorem{corollary}[theorem]{Corollary}
\theoremstyle{definition}
\newtheorem{definition}[theorem]{Definition}
\newtheorem{example}[theorem]{Example}
\theoremstyle{remark}
\newtheorem{remark}[theorem]{Remark}
\newcommand{\CC}{\mathbb C}
\newcommand{\RR}{\mathbb R}
\newcommand{\id}{1}
\newcommand{\HS}{\mathrm{HS}}
\newcommand{\ii}{\mathrm i}
\newcommand{\cK}{\mathcal K}
\newcommand{\cD}{\mathcal D}
\newcommand{\cS}{\mathcal S}
\newcommand{\Nop}{\mathsf N}
\newcommand{\gr}{\varepsilon}
\newcommand{\Eone}{\mathrm E_1}
\newcommand{\mass}{\mathfrak m}
\DeclareMathOperator{\tr}{tr}
\DeclareMathOperator{\str}{str}
\DeclareMathOperator{\Ein}{Ein}
\DeclareMathOperator{\Rea}{Re}
\DeclareMathOperator{\Ima}{Im}
\DeclareMathOperator{\diag}{diag}
\DeclareMathOperator{\spec}{spec}
\DeclareMathOperator{\rk}{rk}
\DeclareMathOperator{\End}{End}
\DeclareMathOperator{\Hom}{Hom}
\newcommand{\norm}[1]{\left\lVert#1\right\rVert}
\newcommand{\inner}[2]{\left\langle#1,#2\right\rangle}


\title[Extending the Kempf--Ness function to complexes]{Extending the
Kempf--Ness function\\ to complexes of quiver representations}
\author{Bernd Johannes Wuebben}
\email{wuebben@gmail.com}
\subjclass[2020]{Primary 16G20, 53D20; Secondary 14F08, 18G80, 53C07}
\keywords{Quiver representation, complex of representations, stability
condition, superconnection, polar decomposition, Kempf--Ness function,
Donaldson functional}
\date{November 18, 2026\\This version: November 18, 2026}

\begin{document}
\begin{abstract}
For every bounded complex of representations of a finite quiver with vertex
charges in the upper half-plane, the superconnection endomorphism of Haiden,
Katzarkov, Kontsevich and Pandit defines a closed one-form on the space of
Hermitian metrics. We construct a primitive by transgression along the
rescaling of the differentials. For nonzero total charge $Z$, rotating the
primitive by $\overline Z/|Z|$ and taking its imaginary part gives a real
potential. On representations this potential is a positive multiple of the
Kempf--Ness function and is nonincreasing along the flow defined by the
argument of the unitary polar factor of the endomorphism. For a two-term
complex in degrees zero and one whose differential surjects onto a term
supported at a loopless vertex $i$, the potential is nonincreasing if the
vertex charge $z_i$ satisfies $\arg z_i\geq\arg Z$ and either
$\arg Z\geq 2\pi/5$ or the differential has rank one. On a three-term
Kronecker quiver complex the potential can increase even when the induced
metric on its stable cohomology is harmonic. Quillen's superconnection
gives the analogous potential for two-term complexes of holomorphic bundles
on a compact Riemann surface. On a single bundle it is a positive multiple
of Donaldson's functional.
\end{abstract}
\maketitle

\section{Introduction}\label{sec:intro}

Haiden, Katzarkov, Kontsevich and Pandit \cite{hkkp26} propose a notion of
K\"ahler geometry on a triangulated category with a stability condition in
the sense of Bridgeland \cite{bri07}. In their framework an object carries a
space of metrics, a complex-valued function $\cS_{\CC}$ on it whose
differential is given by an endomorphism $\Omega$, a phase $\arg\Omega$
defined by a positivity condition, a mass $\tr|\Omega|$, and a flow
$h^{-1}\dot h=-\arg\Omega$ of the metrics, called the minimizing flow in
\cite[\S4.1]{hkkp26}; a metric is \emph{harmonic} if $\arg\Omega$ is a
scalar. For a representation of a quiver with central charges $z_i$ and
Hermitian metrics, $\Omega$ is $z_i$ shifted by the moment map at each
vertex, and the function $\cS_{\CC}$ is explicit \cite[\S5.1]{hkkp26}: after
rotation by the argument of the central charge of the representation, its
imaginary part is a positive multiple of the Kempf--Ness function. For a
complex of representations an endomorphism $\Omega$ is defined in
\cite[\S5.4]{hkkp26} through a superconnection in a graded algebra attached
to the quiver; its trace is the central charge of the complex. No function
$\cS_{\CC}$ is given there for complexes.

For a representation the flow is not the gradient flow of the Kempf--Ness
function, but the Kempf--Ness function does not increase along it
(Proposition~\ref{prop:kempfness}); that is, it is a Lyapunov function for
the flow. The question of this paper is whether the Kempf--Ness function
of a representation extends to complexes. Two properties of the Kempf--Ness
function \cite{kn79,kin94} are meant: its differential is the moment map
shifted by the charges, and it does not increase along the flow. As a
function with differential $\Omega$ it extends to every bounded complex of
representations of every finite quiver (Theorem~\ref{thm:introA}). As a
function that does not increase along the flow it extends to two classes of
complexes (Theorem~\ref{thm:introC}) and not to all
(Theorem~\ref{thm:introD}). Even on a complex with polystable cohomology, the phase potential can be
unbounded below and strictly concave along a geodesic in the space of
metrics (Corollary~\ref{cor:increase}).

In \cite[\S5.4]{hkkp26} the phase of $\Omega$ is written with the functional
calculus, which presupposes that $\Omega$ is normal. For complexes this can
fail (Example~\ref{ex:nonnormal}). In the general framework of
\cite[\S4.2]{hkkp26} the argument is a self-adjoint element $\arg\Omega$ such
that $\Omega(e^{-\sqrt{-1}\arg\Omega}x^*x)\geq0$ for all elements $x$ of the
algebra attached to the metrized object, $\Omega$ being a functional there;
for the functional $x\mapsto\tr(\Omega x)$ on endomorphisms this says that
$\Omega e^{-\ii\arg\Omega}$ is positive semidefinite. We therefore define
$\arg\Omega$ as the Hermitian endomorphism $\Phi$, with spectrum in the
interval $(-\pi n,\pi(1-n))$ in degree $n$, the branch of
\cite[(5.15)]{hkkp26}, for which $\Omega e^{-\ii\Phi}$ is positive definite;
$e^{\ii\Phi}$ is the unitary factor of the polar decomposition of $\Omega$.
For normal $\Omega$ this is the argument of the functional calculus. All
statements about the flow of a complex refer to this definition. We call the
resulting flow the \emph{phase flow}; we do not use the name minimizing flow
because on complexes the mass can increase along it
\cite[Theorem~1.1]{wue26a}. It is the flow of \cite{wue26b,wue26c}, whose
titles keep the name of \cite{hkkp26}.

The phase flow was studied in \cite{wue26a,wue26b} for two-term complexes of
representations of a generalized Kronecker quiver that are quasi-isomorphic
to a representation: the acyclic part decays, and for suitable charges the
flow converges on the cohomology when the cohomology is polystable
\cite[Theorems~1.1 and~1.2]{wue26b}; these two notions are defined in
Section~\ref{sec:setting}. For these complexes a function with
differential $\Omega$ was found by direct computation
\cite[Proposition~7.1]{wue26a}, and for a range of charges its rotated
imaginary part decreases along the flow \cite[Theorem~7.6]{wue26b}.

\subsection{Results}
Let $Q$ be a finite quiver with central charges $z_i$ in the upper
half-plane, and $\mathsf E$ a bounded complex of representations of $Q$ with
spaces $E_i^n$ at the vertex $i$ in degree $n$, differentials $d_i$, arrows
$\rho(a)$ and a Hermitian metric $h_i^n$ on each $E_i^n$
(Section~\ref{sec:setting}). The endomorphism $\Omega$ preserves each
$E_i^n$; we write $\Omega_i^n$ for its restriction to $E_i^n$, and
$\norm\cdot_{\HS}$ for the Hilbert--Schmidt norm of a linear map for the
metrics. An integral formula for $\Omega$ is given in
Proposition~\ref{prop:coefficient}. Its trace is the central charge
$Z=\sum_{i,n}(-1)^nz_i\dim E_i^n$ of the complex.

\begin{theorem}[Primitive]\label{thm:introA}
For every bounded complex of representations of $Q$ with a metric:
\begin{enumerate}
\item[(i)] The one-form $\xi\mapsto\tr(\Omega\,\xi)$ on the space of metrics,
 $\xi=h^{-1}\delta h$, is closed.
\item[(ii)] For $x\in[0,1]$ let $\Omega^{(x)}$ be the endomorphism of the
 complex with differentials $\sqrt x\,d_i$ and the same arrows and metrics.
 Then
 \[
 \cS_{\CC}(h)=\sum_n(-1)^n\Bigl[\sum_iz_i\log\det h_i^n-\sum_a\norm{\rho^n(a)}^2_{\HS}\Bigr]
 +\int_0^1\sum_{i,n}n\,\tr\bigl((\Omega^{(x)})_i^n-(\Omega^{(0)})_i^n\bigr)\,\frac{dx}x
 \]
 satisfies $d\cS_{\CC}(\xi)=\tr(\Omega\,\xi)$. Here $\rho^n(a)$ is the
 restriction of $\rho(a)$ to degree $n$, and the determinants are Gram
 determinants in fixed bases.
\item[(iii)] If $Z$ is not zero, then for any real $\varphi$ with
 $Z=|Z|e^{\ii\varphi}$, along the phase flow
 \[
 \frac d{dt}\Ima\bigl(e^{-\ii\varphi}\cS_{\CC}\bigr)
 =-\sum_{i,n}\tr\Bigl(|\Omega_i^n|\,(\Phi_i^n-\varphi)\sin(\Phi_i^n-\varphi)\Bigr),
 \]
 where $\Omega_i^n=|\Omega_i^n|e^{\ii\Phi_i^n}$ is the polar decomposition
 with the spectrum of $\Phi_i^n$ in the interval
 \mbox{$(-\pi n,\pi(1-n))$}.
\end{enumerate}
\end{theorem}

We call $\Ima(e^{-\ii\varphi}\cS_{\CC})$ the \emph{phase potential}. For a
representation in degree zero, with $\varphi\in(0,\pi)$, it is
$\sin\varphi$ times the Kempf--Ness function
$\sum_a\norm{\rho(a)}^2_{\HS}+\sum_i\theta_i\log\det h_i$ with
$\theta_i=\Ima(e^{-\ii\varphi}z_i)/\sin\varphi$, and each trace in the sum
in (iii) is nonnegative (Proposition~\ref{prop:kempfness}). For complexes
they can have either sign. The last term of $\cS_{\CC}$ is a transgression by
the number operator along the rescaling of the differentials, of the kind
used for Bott--Chern forms of Hermitian complexes \cite{bgs88}, here taken on
the space of metrics and in the algebra attached to the quiver.

Let $K_m$ be the generalized Kronecker quiver, with vertices $0,1$ and $m$
arrows from $0$ to $1$. We write a complex of representations of $K_m$ as a
row of its terms, each term being a representation written as (space at the
vertex $0$ $\to$ space at the vertex $1$); the maps between consecutive terms
are the differentials.

\begin{theorem}[Monotonicity in two classes]\label{thm:introC}
In each of the following cases the phase potential is nonincreasing along the
phase flow.
\begin{enumerate}
\item[(i)] $i_0$ is a vertex of $Q$ without loops; the complex is a vector
 space $E_i^0$ in degree zero at every vertex $i\neq i_0$ and $V\to U$ in
 degrees $0,1$ at $i_0$, with surjective differential and nonzero
 cohomology; and, for $\varphi=\arg Z\in(0,\pi)$, one has
 $\arg z_{i_0}\geq\varphi$ and either $\varphi\geq\frac{\pi}3$ or
 $\dim U=1$.
\item[(ii)] $Q=K_m$ and $\arg z_0\geq\arg z_1$; the complex is
 $(0\to U_1)\to(V\to W)\to(U_0\to0)$ in degrees $-1,0,1$, that is, with $U_1$
 at the vertex $1$ in degree $-1$ and $U_0$ at the vertex $0$ in degree $1$
 (the subscript is the vertex), with injective differential $U_1\to W$,
 surjective differential $V\to U_0$ and nonzero cohomology; and
 $\dim U_0\leq1$, $\dim U_1\leq1$.
\end{enumerate}
\end{theorem}

\begin{theorem}[Failure of monotonicity]\label{thm:introD}
Let $Q=K_m$, $z_0=-1+\ii$ and $z_1=1+\ii$.
\begin{enumerate}
\item[(i)] Let $\mathsf E$ be the orthogonal direct sum of a representation
 with spaces of equal dimension $k\geq1$ at the two vertices, with a harmonic
 metric, and of the complex $\CC\to\CC$ at the vertex $1$ in degrees $0,1$
 with a nonzero differential $f$. Then the phase potential increases along
 the whole trajectory, at the rate $\pi e^{-|f|^2}$, where $|f|$ is the norm
 of $f$ for the metrics at the time considered, and it is not bounded below
 on the metrics of $\mathsf E$.
\item[(ii)] There is a three-term complex as in
 Theorem~\ref{thm:introC}(ii) with $\dim U_0=2$, $\dim U_1=1$ and stable
 cohomology, with a metric at which the phase potential increases along the
 flow.
\end{enumerate}
\end{theorem}

Theorem~\ref{thm:introC} is proved as Theorems~\ref{thm:onevertex}
and~\ref{thm:rankone}, and Theorem~\ref{thm:introD} as
Corollary~\ref{cor:increase} and Theorem~\ref{thm:twosidedfailure}.

\emph{Sharpness.} For $Q=K_m$, $i_0=0$ and cohomology at both vertices,
$\varphi$ lies between $\arg z_1$ and $\arg z_0$, so the condition
$\arg z_0\geq\varphi$ is $\arg z_0\geq\arg z_1$. For $\arg z_0>\arg z_1$,
strict decrease when $\varphi\geq 2\pi/5$ or $\dim U=1$ is
\cite[Theorem~7.6]{wue26b}. The stronger bound $\pi/3$ follows from the elementary trace estimate
in Section~\ref{sec:setting}. Theorem~\ref{thm:introC}(i) adds
arbitrary quivers, the case of equality, and cohomology at a single vertex. In Theorem~\ref{thm:introD}(i) the differential is at
one vertex and $\arg z_1<\varphi$; so the condition
$\arg z_{i_0}\geq\varphi$ of Theorem~\ref{thm:introC}(i) cannot be dropped.
For $\dim U\geq2$ a lower bound on $\varphi$ cannot be dispensed with
either: for $K_2$, the charges $z_0=10^{-6}(-1+\ii)$, $z_1=1+\frac\ii{500}$
and cohomology of dimension vector $(1,1)$, where
$\varphi=0.0020\ldots<\frac{\pi}3$, there is a complex with $\dim U=2$ and a
metric at which the potential increases \cite[Proposition~10.2]{wue26b}. The optimal lower bound on $\varphi$ for
nonnegative total dissipation at every metric remains open. Theorem~\ref{thm:introD}(ii) shows that the
condition on the dimensions in Theorem~\ref{thm:introC}(ii) cannot be
dropped; for $z_0=-1+\ii$, $z_1=1+\ii$ and cohomology with equal dimensions
at the two vertices, the pairs $(\dim U_0,\dim U_1)$ for which the phase
potential is nonincreasing at all metrics are determined in
Corollary~\ref{cor:classification}.

\emph{Method.} The closedness in Theorem~\ref{thm:introA} needs only two odd
derivations of a graded algebra with the relations of
Lemma~\ref{lem:abstract}, a functional that vanishes on
supercommutators and on the images of the derivations, and a formula for the
variation of the curvature element under a change of metric; the same lemma
serves the algebra of the quiver
and, for bundles, the algebra of endomorphism-valued forms, where it is the
classical argument by which Bott--Chern forms do not depend, modulo
$\partial$- and $\bar\partial$-exact forms, on the path of metrics. Theorem~\ref{thm:introC}(i) uses the trace estimate
proved after Theorem~\ref{thm:imported}, which strengthens
\cite[Theorem~3.1]{wue26b}. The entries of the Schur multiplier are the
second divided differences of $e^{-x}$ at $\ell_i,0,\ell_k$, where the
$\ell_i$ are the eigenvalues of the differential times its adjoint. In the three-term complexes the point $0$ is
replaced by the eigenvalues of the other differential times its adjoint, and
no such inequality holds uniformly in these eigenvalues
(Remark~\ref{rem:interpolation}); this is the source of the example in
Theorem~\ref{thm:introD}(ii). The numerical values in that example are
enclosures obtained with interval arithmetic \cite{joh17}.

\begin{theorem}[Duality]\label{thm:introB}
Let $\mathsf E^\vee$ be the dual complex over the opposite quiver, with
$(E^\vee)_i^n=(E_i^{-n})^*$, the transposed maps, the dual metrics and the
charges $-\overline{z_i}$. Then $\Omega^\vee=-(\Omega^\dagger)^{\mathsf T}$,
the phases satisfy $\Phi^\vee=\pi-\Phi^{\mathsf T}$, the masses agree, and
$h\mapsto e^{-\pi t}(h^{-1})^{\mathsf T}$ maps trajectories of the phase flow
of $\mathsf E$ to trajectories of the phase flow of $\mathsf E^\vee$. If
$Z(\mathsf E)\neq0$ and the argument of $Z(\mathsf E^\vee)$ is taken as
$\pi-\varphi$, this map preserves the phase potential up to a constant.
\end{theorem}

Duality is used as a tool: it halves the cases in
Theorem~\ref{thm:onevertex}, and it carries statements about complexes with a
surjective differential in degrees $0,1$ over to complexes with an injective
differential in degrees $-1,0$ (Remark~\ref{rem:injective}).

\subsection{Bundles}
The last section gives the corresponding formulation for a two-term complex
$E^0\to E^1$ of holomorphic vector bundles with Hermitian metrics on a
compact Riemann surface, with $-\frac{\ii}{2\pi}\Lambda F_h$ in the place of
the moment map, following the dictionary of \cite[Remark~5.14]{hkkp26} and
using Quillen's superconnection \cite{qui85}. We record the formal
properties, which for the closedness follow the theory of Bott--Chern forms
\cite{bc65,don85,bgs88}: the one-form defined by $\Omega$ on the space of
metrics is closed, with a primitive given by the same transgression; for a
single bundle the phase potential is a positive multiple of Donaldson's
functional \cite{don85}; the derivative formula of
Theorem~\ref{thm:introA}(iii) holds along smooth solutions; the
linearization is parabolic in the sense of Petrovskii
for the convention $\Delta_\omega=-d^*d$; and for a section $s$ of a
line bundle the sum $(\Omega^0+\Omega^1)\,\omega$ of the two blocks of
$\Omega$, times the K\"ahler form, is a representative of the Euler class of
the kind constructed by Mathai and Quillen \cite{mq86}, and the integrand of
the phase potential decreases at every point of the surface
(Propositions~\ref{prop:bundleprimitive}--\ref{prop:section}).

\subsection{Open questions}
Convergence of the phase flow on the cohomology and decay of the acyclic
part are not treated here, except for the contractible summands of
Proposition~\ref{prop:contractible}. For two-term complexes over $K_m$ with
a surjective differential at the vertex $0$, the decay is proved in
\cite{wue26b} for all charges and the convergence for strictly ordered
charges and polystable cohomology. For the corresponding complexes on
acyclic quivers, \cite{wue26c} proves decay of the exact part and
exponential convergence of the rescaled metrics on polystable
Harder--Narasimhan quotients of the cohomology.
Remark~\ref{rem:injective} carries these results to injective differentials. For the
three-term complexes of Theorem~\ref{thm:introC}(ii) both are open. Beyond
Theorems~\ref{thm:onevertex} and~\ref{thm:rankone} and
Corollaries~\ref{cor:increase} and~\ref{cor:classification} we do not know
when the phase potential is monotone. For bundles, existence of solutions of
the flow is not proved, and neither is the convergence along the flow, for a
section of a line bundle, of the form $(\Omega^0+\Omega^1)\,\omega$ to the
current of integration over the zeros of the section.

\subsection{Organization}
Section~\ref{sec:setting} recalls the definition of \cite{hkkp26}, gives the
integral formula for $\Omega$, fixes the terminology, and sharpens one of the two
matrix inequalities recalled from \cite{wue26b}. Section~\ref{sec:primitive}
proves Theorem~\ref{thm:introA}(i),(ii), and Section~\ref{sec:rate} proves
(iii) and Theorem~\ref{thm:introD}(i). Section~\ref{sec:duality} proves
Theorem~\ref{thm:introB}. Section~\ref{sec:monotone} proves
Theorems~\ref{thm:introC} and~\ref{thm:introD}(ii).
Section~\ref{sec:bundles} is on bundles.

\section{The endomorphism \texorpdfstring{$\Omega$}{Omega} of a complex of representations}\label{sec:setting}

\subsection{Complexes of representations and metrics}
Let $Q$ be a finite quiver with set of vertices $Q_0$ and set of arrows
$Q_1$; we write $a:i\to j$ for an arrow from $i$ to $j$. Loops and oriented
cycles are allowed. We write $K_m$ for the generalized Kronecker quiver with
vertices $0,1$ and $m\geq1$ arrows $a_1,\dots,a_m$ from $0$ to $1$. Fix
\emph{central charges} $z_i$, $i\in Q_0$, in the open upper half-plane, and
put $\alpha_i=\arg z_i\in(0,\pi)$; they define the central charge of a stability
condition on the bounded derived category of finite-dimensional
representations of $Q$ whose heart is the category of representations
\cite[Propositions~2.4 and~5.3]{bri07}. Phases in this paper are arguments, in radians; in the
normalization of \cite{bri07} they are divided by $\pi$, and a shift by one
degree changes them by $\pi$. A \emph{complex of representations}
$\mathsf E$ consists of finite-dimensional complex vector spaces $E_i^n$
($i\in Q_0$, $n\in\mathbb Z$, almost all zero), differentials
$d_i:E_i^n\to E_i^{n+1}$ with $d_i^2=0$, and for each arrow $a:i\to j$ a chain
map $\rho(a):E_i^\bullet\to E_j^\bullet$ of degree zero. Its central charge is
\[
 Z(\mathsf E)=\sum_{i,n}(-1)^nz_i\dim E_i^n .
\]
A \emph{metric} $h$ on $\mathsf E$ is a Hermitian metric $h_i^n$ on each
$E_i^n$; a complex with a metric is called metrized. The adjoint of a linear
map for the metrics is written ${}^\dagger$, inner products are
conjugate-linear in the first argument, and
$\norm T_{\HS}=(\tr T^\dagger T)^{1/2}$ is the Hilbert--Schmidt norm. For an
endomorphism $T$ of a Hermitian space, $\frac12(T+T^\dagger)$ is its
Hermitian part and $\Ima T=\frac1{2\ii}(T-T^\dagger)$ its imaginary part. We
write $E_i=\bigoplus_nE_i^n$
and $E=\bigoplus_iE_i$. The \emph{number operator} $\Nop$ is the endomorphism
of $E$ equal to $n$ on $E_i^n$, and $\gr=(-1)^{\Nop}$ is the grading. A
tangent vector to the space of metrics at $h$ is an endomorphism
$\xi=h^{-1}\delta h$ of $E$ that preserves each $E_i^n$ and is Hermitian for
$h$. In the algebra $\mathcal A(Q)$ below we write $\sqrt{-1}$ for the
imaginary unit, as in \cite{hkkp26}; elsewhere it is $\ii$, and the italic
$i$ is a vertex.

\subsection{The definition of Haiden, Katzarkov, Kontsevich and Pandit}
The graded $*$-algebra $\mathcal A(Q)$ of \cite[Definition~5.13]{hkkp26} has
the basis $e_i,f_i$ ($i\in Q_0$) in degree $0$, $a\in Q_1$ in degree $1$ and
$a^*$ in degree $-1$, and the nonzero products of basis elements are
\begin{gather*}
 e_i^2=e_i,\qquad e_if_i=f_ie_i=f_i,\qquad e_ja=ae_i=a,\qquad e_ia^*=a^*e_j=a^*,\\
 aa^*=\sqrt{-1}\,f_j,\qquad a^*a=-\sqrt{-1}\,f_i
\end{gather*}
for $a:i\to j$. In the graded tensor product
$\mathcal A(Q)\otimes\End(E)$ the product is
$(x\otimes T)(y\otimes S)=(-1)^{|T||y|}xy\otimes TS$. The involution of
$\mathcal A(Q)$ fixes the $e_i$ and $f_i$ and exchanges $a$ and $a^*$; on
$\End(E)$ it is $T^*=T^\dagger$ for even $T$ and $T^*=\sqrt{-1}\,T^\dagger$
for odd $T$, the adjoints being those of the metric. Let
$\mathcal B\subset\mathcal A(Q)\otimes\End(E)$ be the subalgebra
\[
 \mathcal B=\bigoplus_i\bigl(e_i\otimes\End(E_i)\oplus f_i\otimes\End(E_i)\bigr)
 \oplus\bigoplus_{a:i\to j}\bigl(a\otimes\Hom(E_i,E_j)\oplus a^*\otimes\Hom(E_j,E_i)\bigr),
\]
with unit $\sum_ie_i\otimes\id_{E_i}$. Put
\[
 A=\sum_ie_i\otimes d_i+\sum_aa\otimes\rho(a),\qquad
 F=(A-A^*)^2+\sqrt{-1}\sum_if_i\otimes z_i\id_{E_i} .
\]
Here $A^*=\sqrt{-1}\sum_ie_i\otimes d_i^\dagger+\sum_aa^*\otimes\rho(a)^\dagger$,
and this is the only use of the involution. These are elements of
$\mathcal B$. Following \cite[\S5.4]{hkkp26} we put
\begin{equation}\label{eq:Omegadef}
 \Omega:=\gr\,\bigl[\exp(-\sqrt{-1}\,F)\bigr]_f\in\End(E),
\end{equation}
where the exponential is taken in $\mathcal B$ and, for
$X=\sum_i(e_i\otimes X_{e_i}+f_i\otimes X_{f_i})+\dots\in\mathcal B$, the
endomorphism $[X]_f=\bigoplus_iX_{f_i}$ is the component along the $f_i$.
With
\begin{equation}\label{eq:DeltaMGamma}
 \Delta_i=d_id_i^\dagger+d_i^\dagger d_i,\qquad
 M_a=d_j^\dagger\rho(a)-\rho(a)d_i^\dagger,\qquad
 \Gamma_i=z_i+\sum_{a:i\to\cdot}\rho(a)^\dagger\rho(a)-\sum_{a:\cdot\to i}\rho(a)\rho(a)^\dagger ,
\end{equation}
where $M_a:E_i^\bullet\to E_j^{\bullet-1}$ for $a:i\to j$, one has
\cite[display after Remark~5.15]{hkkp26}
\begin{equation}\label{eq:F}
 -\sqrt{-1}\,F=-\sum_ie_i\otimes\Delta_i+\sum_aa\otimes M_a
 +\sqrt{-1}\sum_aa^*\otimes M_a^\dagger+\sum_if_i\otimes\Gamma_i .
\end{equation}
For a representation, that is, a complex in degree zero, $\Omega=\Gamma$:
at each vertex $\Omega$ is $z_i$ shifted by the moment map of the action of
the unitary group of $E_i$.

In \cite[Remark~5.14]{hkkp26} the algebra $\mathcal A(Q)$ is compared with
the differential forms on a Riemann surface: the $e_i$ correspond to
functions, $a$ and $a^*$ to forms of type $(1,0)$ and $(0,1)$, and the $f_i$
to two-forms. In this dictionary $A-A^*$ plays the role of a superconnection
in the sense of Quillen \cite{qui85} on the graded space $E$, $F$ that of
its curvature shifted by the central term, $[\,\cdot\,]_f$ that of the
component of top degree, and the functional $\tau$ of
Section~\ref{sec:primitive} that of the integral of the supertrace of this
component. Section~\ref{sec:bundles} carries the dictionary out for bundles.

\subsection{An integral formula}
For Hermitian $\Delta\geq0$ and an operator $G$ put
$K_\Delta(G)=\int_0^1e^{-(1-x)\Delta}Ge^{-x\Delta}dx$, and for operators
$X:E'\to E''$, $Y:E''\to E'$ and Hermitian $\Delta',\Delta''\geq0$ on $E'$,
$E''$,
\[
 \mathcal J_{\Delta'',\Delta'}(X,Y)
 =\int_{\substack{x,y\geq0\\x+y\leq1}}e^{-x\Delta''}X\,e^{-(1-x-y)\Delta'}\,Y\,e^{-y\Delta''}\,dx\,dy .
\]
In eigenbases of $\Delta$, $\Delta''$ and $\Delta'$, with eigenvalues
$\ell_i$, $\ell''_i$ and $\ell'_\mu$, these are given by divided differences
of the function $\mathsf f(x)=e^{-x}$, written $\mathsf f[\,\cdot\,]$, by the
Hermite--Genocchi formula:
\[
 K_\Delta(G)_{ik}=-\mathsf f[\ell_i,\ell_k]\,G_{ik},\qquad
 \mathcal J_{\Delta'',\Delta'}(X,Y)_{ik}=\sum_\mu\mathsf f[\ell''_i,\ell'_\mu,\ell''_k]\,X_{i\mu}Y_{\mu k},
\]
with $-\mathsf f[a,b]>0$ and $\mathsf f[a,b,c]>0$, by the mean value form of
divided differences. The expansion
$e^{X+Y}=e^X+\int_0^1e^{(1-x)X}Ye^{xX}dx+\dots$ is called the Duhamel
expansion below.

\begin{proposition}\label{prop:coefficient}
The endomorphism $\Omega$ preserves each $E_i^n$, and for every vertex $i$
\[
 \gr\,\Omega\big|_{E_i}=K_{\Delta_i}(\Gamma_i)
 +\sum_{a:k\to i}\mathcal J_{\Delta_i,\Delta_k}(M_a,M_a^\dagger)
 -\sum_{a:i\to j}\mathcal J_{\Delta_i,\Delta_j}(M_a^\dagger,M_a).
\]
\end{proposition}

\begin{proof}
Write the right side of \eqref{eq:F} as $X_e+X_a+X_{a^*}+X_f$. In
$\mathcal A(Q)$ every product of three elements of $\{a,a^*,f_i\}$ vanishes,
and so do the products $f_if_k$, $f_ia$, $af_i$, $f_ia^*$, $a^*f_i$, $aa'$,
$a^*a'^*$, and $aa'^*$, $a^*a'$ for $a\neq a'$. So in the Duhamel expansion of
$\exp(X_e+Y)$, $Y=X_a+X_{a^*}+X_f$, the component along the $f_i$ receives
only the terms with one factor $X_f$, or with one factor $X_a$ and one
factor $X_{a^*}$ of the same arrow, in either order. The element $X_e$ is
even and the $e_i$ act as units, so
$e^{xX_e}=\sum_ie_i\otimes e^{-x\Delta_i}$ and no signs arise from it. By the
sign rule, for $a:i\to j$,
\begin{align*}
 (a\otimes M_a)(\sqrt{-1}\,a^*\otimes M_a^\dagger)
 &=\sqrt{-1}\,(-1)^{|M_a||a^*|}aa^*\otimes M_aM_a^\dagger=f_j\otimes M_aM_a^\dagger,\\
 (\sqrt{-1}\,a^*\otimes M_a^\dagger)(a\otimes M_a)
 &=\sqrt{-1}\,(-1)^{|M_a^\dagger||a|}a^*a\otimes M_a^\dagger M_a=-f_i\otimes M_a^\dagger M_a .
\end{align*}
The second-order term of the expansion is
$\int e^{xX_e}Ye^{(1-x-y)X_e}Ye^{yX_e}\,dx\,dy$ over the simplex; inserting the
heat factors gives the three terms. Each term preserves the degree, because
$\Delta_i,\Gamma_i$ have degree $0$, $M_a$ has degree $-1$ and $M_a^\dagger$
degree $+1$.
\end{proof}

The formula is the expansion to second order of the display for
$-\sqrt{-1}F$ in \cite[\S5.4]{hkkp26}; no higher terms occur.

\begin{proposition}\label{prop:properties}
For every metrized complex:
\begin{enumerate}
\item[(a)] All terms of Proposition~\ref{prop:coefficient} other than the
 summand $K_{\Delta_i}(z_i)=z_ie^{-\Delta_i}$ of $K_{\Delta_i}(\Gamma_i)$ are
 Hermitian. Hence, on each $E_i^n$,
 $\Ima\bigl((-1)^n\Omega\bigr)=(\Ima z_i)\,e^{-\Delta_i}>0$.
\item[(b)] $\tr\Omega=Z(\mathsf E)$ \cite[Lemma~5.17]{hkkp26}.
\item[(c)] On each $E_i^n$ there is a unique Hermitian $\Phi_i^n$ with
 spectrum in $(-\pi n,\pi(1-n))$ such that
 $\Omega|_{E_i^n}\,e^{-\ii\Phi_i^n}$ is positive definite.
\item[(d)] The flow $h^{-1}\dot h=-\Phi$, $\Phi=\bigoplus\Phi_i^n$, has a
 unique solution for all $t\geq0$ from every initial metric.
\end{enumerate}
\end{proposition}

\begin{proof}
(a) $K_\Delta(G)^\dagger=K_\Delta(G^\dagger)$ by the substitution
$x\mapsto1-x$, and $\mathcal J(X,X^\dagger)^\dagger=\mathcal J(X,X^\dagger)$ by
exchanging $x$ and $y$. (b) This is \cite[Lemma~5.17]{hkkp26}; the argument
is as follows. With $\tau(X)=\tr(\gr\,[X]_f)$ one has
$\tr\Omega=\tau(e^{-\sqrt{-1}F})$. Replace all structure maps by their
multiples by $u\in[0,1]$; then $F$ becomes $-u^2[A,A^*]+F_z$ with
$F_z=\sqrt{-1}\sum_if_i\otimes z_i$, the element $A$ commutes with it, and the
derivative of $\tau(e^{-\sqrt{-1}F})$ in $u$ is
$2\sqrt{-1}\,u\,\tau\bigl([A,A^*e^{-\sqrt{-1}F}]\bigr)=0$, because $\tau$
vanishes on supercommutators of elements of $\mathcal B$
(Lemma~\ref{lem:tau} below). At $u=0$,
$e^{-\sqrt{-1}F_z}=\sum_i(e_i+f_i\otimes z_i)$ and
$\tau=\sum_{i,n}(-1)^nz_i\dim E_i^n$. (c) Let $T$ be a matrix with $\Ima T$ positive
definite. Then $T$ is invertible, since $Tv=0$ gives
$\inner v{(\Ima T)v}=\Ima\inner v{Tv}=0$. Let $T=PU$ be its polar
decomposition, with $P=(TT^\dagger)^{1/2}>0$ and $U$ unitary. If
$Uv=\lambda v$, then $\inner v{(\Ima T)v}=(\Ima\lambda)\inner v{Pv}$, so
$\Ima\lambda>0$. Hence $U=e^{\ii\Psi}$ for a unique Hermitian $\Psi$ with
spectrum in $(0,\pi)$, and $Te^{-\ii\Psi}=P>0$; conversely such a $\Psi$
gives a polar decomposition, which is unique. Apply this to
$(-1)^n\Omega|_{E_i^n}$ and put $\Phi_i^n=\Psi-\pi n$. (d) The vector field is real-analytic on the space of
metrics, and $\norm{h^{-1}\dot h}\leq\pi(1+\max|n|)$, so solutions cannot
leave compact sets in finite time.
\end{proof}

The computation in (a) is that of \cite[proof of Lemma~5.16]{hkkp26}. That
lemma also states that $\Omega$ is normal. By (a), a block $\Omega|_{E_i^n}$
is normal if and only if its Hermitian part commutes with $e^{-\Delta_i}$,
and this does not hold for all complexes:

\begin{example}[$\Omega$ need not be normal]\label{ex:nonnormal}
Let $Q=K_1$, with charges $z_0=-1+\ii$, $z_1=1+\ii$, and consider the complex
with the spaces $V=\CC^2$ in degree $0$ and $U_0=\CC^2$ in degree $1$ at the
vertex $0$, the spaces $U_1=\CC$ in degree $-1$ and $W=\CC$ in degree $0$ at
the vertex $1$, the standard metrics, the differentials
$d_0=\diag(1,4):V\to U_0$ and $d_1=3:U_1\to W$, and the arrow
$\rho(a_1)=(\rho_1,\rho_2)=(4,100):V\to W$. On $V$ one has
$d_0^\dagger d_0=\diag(\ell_1,\ell_2)$ with $(\ell_1,\ell_2)=(1,16)$ and
$\Gamma_0=z_0+\rho(a_1)^\dagger\rho(a_1)$; the operator $M_{a_1}$ of
\eqref{eq:DeltaMGamma} restricts to $3\rho(a_1):V\to U_1$, and $\Delta_1=9$
on $U_1$. By Proposition~\ref{prop:coefficient} (see
Proposition~\ref{prop:twosidedcoefficient} for the general formula), the
block of $\Omega$ on $V$ is
\[
 (\Omega_V)_{ik}=z_0e^{-\ell_i}\delta_{ik}
 +\rho_i\rho_k\bigl(-\mathsf f[\ell_i,\ell_k]-9\,\mathsf f[\ell_i,9,\ell_k]\bigr),
\]
where the factor $9$ is $|d_1|^2$ and the point $9$ is the eigenvalue of
$d_1^\dagger d_1$. Its imaginary part $\diag(e^{-1},e^{-16})$ has distinct
eigenvalues, and its Hermitian part has the off-diagonal entry
$400\,(-\mathsf f[1,16]-9\,\mathsf f[1,9,16])=-1.2183\ldots\neq0$; so the two
do not commute. This complex is used again in
Theorem~\ref{thm:twosidedfailure}. For two-term complexes over $K_m$ see
\cite[Example~2.3]{wue26a}.
\end{example}

\subsection{Terminology}

\begin{definition}\label{def:phase}
Let $\mathsf E$ be a metrized complex.
\begin{enumerate}
\item[(i)] The \emph{blocks} of $\Omega$ are the endomorphisms
 $\Omega_i^n=\Omega|_{E_i^n}$. The \emph{phase} is
 $\Phi=\bigoplus\Phi_i^n$ with $\Phi_i^n$ as in
 Proposition~\ref{prop:properties}(c), and
 $|\Omega_i^n|:=\Omega_i^ne^{-\ii\Phi_i^n}=(\Omega_i^n\Omega_i^{n\dagger})^{1/2}$.
 The word phase is also used, as usual, for the argument of a central
 charge.
\item[(ii)] The \emph{phase flow} is the flow $h^{-1}\dot h=-\Phi$ of
 Proposition~\ref{prop:properties}(d). The \emph{mass} is
 $\mass(h)=\sum_{i,n}\tr|\Omega_i^n|$, and a metric is \emph{harmonic} if
 $\Phi$ is a scalar.
\item[(iii)] Let $i$ be a vertex. The complex of vector spaces $E_i^\bullet$
 splits orthogonally as $\ker\Delta_i\oplus(\ker\Delta_i)^\perp$; the
 differential vanishes on the first summand, which maps isomorphically onto
 the cohomology, and the second summand is an acyclic subcomplex. We call
 $(\ker\Delta_i)^\perp$ the \emph{acyclic part of $\mathsf E$ at $i$} and the
 rank of $d_i$ its \emph{rank}. The metrics \emph{induced on the cohomology}
 are the restrictions of the metrics to the spaces $\ker\Delta_i$.
\item[(iv)] A \emph{contractible summand} is a direct summand of
 $\mathsf E$, as a complex of representations, that is contractible. Those
 used below are complexes $C^n\to C^{n+1}$ at one vertex, with invertible
 differential, on which no arrow acts.
\item[(v)] Suppose that the cohomology of $\mathsf E$ is a representation in
 degree zero, with $Z(\mathsf E)=|Z|e^{\ii\varphi}$, $\varphi\in(0,\pi)$.
 Along a trajectory of the phase flow, the acyclic part \emph{decays} if the
 least nonzero eigenvalue of each $\Delta_i$ tends to infinity and $\Omega$
 minus the endomorphism of the cohomology representation with its induced
 metrics, extended by zero to the acyclic parts, tends to zero in the trace
 norm $\tr(T^\dagger T)^{1/2}$ for the metrics $h(t)$ (this is called
 cancellation of the acyclic part
 in \cite{wue26b} and decay of the exact part in \cite{wue26c}); and the
 flow
 \emph{converges on the cohomology} if the metrics induced on the
 cohomology, multiplied by $e^{\varphi t}$, converge to a harmonic metric of
 the cohomology representation.
\end{enumerate}
\end{definition}

\begin{remark}\label{rem:phase}
(1) The phase of Definition~\ref{def:phase} satisfies the positivity
condition of \cite[\S4.2]{hkkp26} recalled in the Introduction, with the
branches of \cite[(5.15)]{hkkp26}, and the phase flow is the minimizing flow
of \cite[\S4.1]{hkkp26} for this phase; for normal $\Omega$ the phase is the
argument in the sense of the functional calculus used in
\cite[\S5.4]{hkkp26}.

(2) Since the intervals $(-\pi n,\pi(1-n))$ are pairwise disjoint, a harmonic
metric can exist only on a complex concentrated in one degree $n$, that is,
on a shifted representation; then $\Phi$ is the scalar $\arg Z(\mathsf E)$
taken in $(-\pi n,\pi(1-n))$, the metric is a zero of the moment map shifted
by the stability parameter (Proposition~\ref{prop:kempfness}), and one
exists if and only if the representation is polystable
\cite[Theorem~5.1]{hkkp26}, \cite[Proposition~6.5]{kin94}. For a complex with
terms in several degrees, `harmonic' is used below only for the metrics
induced on the cohomology.

(3) The acyclic part at a vertex is in general not preserved by the arrows,
and an acyclic direct summand need not be contractible
(Theorem~\ref{thm:twosidedfailure}).
\end{remark}

Changing the metric at fixed structure maps is the same as moving the
structure maps in their orbit under $\prod\mathrm{GL}(E_i^n)$ at a fixed
metric. Precisely, if $h=\overline g^{\mathsf T}g$ in fixed bases, the matrix
of a structure map
$T\in\{d_i,\rho(a)\}$ in the orthonormal frames given by $g$ is
$g_{\mathrm{target}}Tg_{\mathrm{source}}^{-1}$; along the phase flow one can
choose $g$ with $\dot gg^{-1}=-\frac12\Phi$, where $\Phi$ now denotes the
Hermitian matrix of the phase in these frames, and then
\begin{equation}\label{eq:frameflow}
 \dot T=\tfrac12\bigl(T\,\Phi_{\mathrm{source}}-\Phi_{\mathrm{target}}\,T\bigr)
\end{equation}
for the matrices in orthonormal frames, where $\Phi_{\mathrm{source}}$,
$\Phi_{\mathrm{target}}$ are the phases on the source and the target of $T$.
The endomorphism $\Omega$, the phases and the mass depend only on the
matrices of the structure maps in orthonormal frames, up to unitary changes
of the frames.

For $Q=K_m$ the charges are called \emph{ordered} if
$\alpha_0\geq\alpha_1$ and \emph{strictly ordered} if $\alpha_0>\alpha_1$;
that is, the simple representation at the source vertex has phase at least,
respectively larger than, that of the simple representation at the target
vertex. We
write a complex of representations of $K_m$ as a row of its terms, each term
being the representation $(E_0^n\to E_1^n)$ with the space at the vertex $0$
on the left; the maps between consecutive terms are the differentials
$(d_0,d_1)$.

\subsection{Two matrix inequalities}
We recall two matrix inequalities of \cite{wue26b} and strengthen the first. For Hermitian matrices $L>0$
and $M$ put
\[
 J_L(M)=\int_{\substack{x,y\geq0\\x+y\leq1}}e^{-xL}Me^{-yL}\,dx\,dy ,
\]
so that $J_L(XY)=\mathcal J_{L,0}(X,Y)$; in an eigenbasis of $L$ with
eigenvalues $\ell_i$, $J_L(M)_{ik}=\mathsf f[\ell_i,0,\ell_k]\,M_{ik}$. We
write $M_\pm\geq0$ for the positive and negative parts of a Hermitian matrix
$M=M_+-M_-$.

\begin{theorem}[{\cite[Theorems~3.1 and~4.1]{wue26b}}]\label{thm:imported}
The following hold.
\begin{enumerate}
\item[(a)] For Hermitian matrices $L>0$ and $M\geq0$,
 \begin{equation}\label{eq:kernel}
 \tr J_L(M)_-\leq\tfrac25\tr J_L(M).
 \end{equation}
\item[(b)] Let $T(s)$ be a smooth family of matrices with constant positive
 definite imaginary part, and $\Psi(s)$ the Hermitian matrices with spectrum
 in $(0,\pi)$ and $T e^{-\ii\Psi}>0$. Then $\tr(\dot T\,\dot\Psi)<0$ wherever
 $\dot T\neq0$.
\end{enumerate}
\end{theorem}

For $1\times1$ matrices, (a) is empty because $J_L(M)\geq0$, and (b) says
that $\arg(x+\ii y)$ decreases in $x$ for fixed $y>0$. In \cite{wue26b}, (a)
is reduced to an inequality in one real variable, proved with the help of
interval arithmetic, and (b) is proved by expressing
$\tr(\dot T\,\dot\Psi)$ through a Sylvester equation as a quadratic form whose
kernels are positive semidefinite by an integral representation
\cite[Section~4]{wue26b}.

We shall use the stronger estimate
$\tr J_L(M)_-\leq\frac13\tr J_L(M)$, which admits a direct proof.
In an eigenbasis of $L$, write $J_L(M)_{ik}=M_{ik}j(\ell_i,\ell_k)$,
where $j(a,b)$ is the integral of $e^{-ax-by}$ over $x,y\geq0$, $x+y\leq1$.
For $s\in I=[0,1/2]$ put $F_a(s)=\int_0^s e^{-ax}\,dx$,
$A_a=F_a(1/2)$ and $V_a(s)=e^{-a(1-s)}$. Splitting the triangle at
$x=1/2$ and $y=1/2$ gives
$j(a,b)=A_aA_b+\int_I(F_aV_b+V_aF_b)\,ds$. Thus both $Q$ and $j+Q$ are
positive semidefinite kernels, where
\[
 Q(a,b)=\tfrac23\int_I(F_a-\tfrac34V_a)(F_b-\tfrac34V_b)\,ds .
\]
Indeed, $j+Q$ is $A_aA_b$ plus the same integral with both minus signs
replaced by plus signs. On the diagonal suppress the subscript $a$. Since
$sV\leq F\leq A$ and $V^2$ is increasing,
$\int_I FV\geq\int_I sV^2\geq\tfrac14\int_I V^2$. Expanding the square gives
\[
 Q(a,a)\leq\tfrac13 A^2+\tfrac12\int_I FV\,ds
 \leq\tfrac13 j(a,a).
\]
By the Schur product theorem, the matrices with entries
$M_{ik}(j+Q)(\ell_i,\ell_k)$ and $M_{ik}Q(\ell_i,\ell_k)$ are positive
semidefinite. Taking the trace on the negative spectral subspace of their
difference bounds $\tr J_L(M)_-$ by $\sum_i M_{ii}Q(\ell_i,\ell_i)$,
which is at most $\frac13\tr J_L(M)$.
This estimate enters Theorem~\ref{thm:onevertex}(b),(c) in rank at least
two and hence the one-sided cases of Corollary~\ref{cor:classification}.
Proposition~\ref{prop:parabolic} uses Theorem~\ref{thm:imported}(b).

\section{Closedness and a primitive}\label{sec:primitive}

Consider the complex-valued one-form
\[
 \vartheta_h(\xi):=\tr(\Omega\,\xi)
\]
on the space of metrics; here $\xi=h^{-1}\delta h$ is a tangent vector at
$h$. For representations $\Omega=\Gamma$, and
$\vartheta$ is the differential of
$\sum_iz_i\log\det h_i-\sum_a\norm{\rho(a)}_{\HS}^2$ \cite[\S5.1]{hkkp26}. For
two-term complexes over $K_m$ with a surjective differential at the vertex
$0$, a function with differential $\vartheta$ was found by direct computation
in \cite[Proposition~7.1]{wue26a}. Here closedness is proved for all
complexes by a computation in the algebra $\mathcal B$ of
Section~\ref{sec:setting}.

The algebra $\mathcal B$ contains $A$, $A^*$, $F$ and the elements
$\sum_ie_i\otimes\xi_i$, with which we identify tangent vectors $\xi$. For
$X\in\mathcal B$ let $\tau(X)=\tr\bigl(\gr\,[X]_f\bigr)$, so that
$\vartheta_h(\xi)=\tau(e^{-\sqrt{-1}F}\xi)$. All brackets below are
supercommutators, $[X,Y]=XY-(-1)^{|X||Y|}YX$.

\begin{lemma}\label{lem:tau}
$\tau$ vanishes on supercommutators of elements of $\mathcal B$.
\end{lemma}

\begin{proof}
The products of basis elements with a component along some $f_i$ are those of
the types $(e_i,f_i)$, $(f_i,e_i)$, $(a,a^*)$ and $(a^*,a)$. For the first
two the claim is that the supertrace on $E_i$ vanishes on supercommutators.
For $a:i\to j$, $T\in\Hom(E_i,E_j)$ and
$S\in\Hom(E_j,E_i)$ homogeneous with $|T|+|S|=0$,
\[
 \tau\bigl([a\otimes T,a^*\otimes S]\bigr)
 =\sqrt{-1}\bigl((-1)^{|T|}\str_{E_j}(TS)-\str_{E_i}(ST)\bigr)=0,
\]
because $\str_{E_j}(TS)=(-1)^{|T|}\str_{E_i}(ST)$; here $\str=\tr(\gr\,\cdot\,)$.
\end{proof}

The restriction to $\mathcal B$ matters. On the whole algebra
$\mathcal A(Q)\otimes\End(E)$, with the projection onto
$\bigoplus_if_i\otimes\End(E_i)$ in the place of $[\,\cdot\,]_f$, the functional
does not vanish on all supercommutators, for instance on
$[e_i\otimes T,f_i\otimes S]$ with $T:E_j\to E_i$, $S:E_i\to E_j$, $i\neq j$.

The closedness of $\vartheta$ uses only the structure described in the next
lemma. For the endomorphism-valued forms of Section~\ref{sec:bundles} the
lemma is the classical computation by which the one-form
$\xi\mapsto\str(\xi\,e^{-\mathfrak F})$ on the space of metrics is closed
modulo $\partial$- and $\bar\partial$-exact forms, and by which Bott--Chern
forms do not depend, modulo such forms, on the path of metrics
\cite{bc65,don85,bgs88}. We
isolate the identities used in that argument because they also hold in the
algebra $\mathcal B$.

\begin{lemma}\label{lem:abstract}
Let $\mathcal C$ be a $\mathbb Z/2$-graded associative algebra over $\CC$
with a linear functional $\tau$ that vanishes on supercommutators. For each
point $h$ of a manifold $\mathcal M$ let $\mathfrak d$, $\mathfrak d^\star$
be odd derivations of $\mathcal C$ and $F\in\mathcal C$ an even element,
depending smoothly on $h$, such that
\[
 \tau\circ\mathfrak d=0=\tau\circ\mathfrak d^\star,\qquad
 \mathfrak dF=0=\mathfrak d^\star F,\qquad
 \mathfrak d\mathfrak d^\star+\mathfrak d^\star\mathfrak d=-[F,\,\cdot\,].
\]
Suppose that to each tangent vector of $\mathcal M$ at $h$ an even element
$\xi\in\mathcal C$ is assigned, linearly, such that the derivative of $F$
in this direction is $\delta_\xi F=-\mathfrak d\mathfrak d^\star\xi$, and
that for a two-parameter family $h(u,v)$, with the elements $\xi$, $\eta$
assigned to $\partial_uh$, $\partial_vh$, one has
$\partial_v\xi-\partial_u\eta=[\xi,\eta]$. Suppose finally that
$\mathsf e:=e^{-\sqrt{-1}F}$ and the elements
\[
 \Xi_Y:=\int_0^1e^{-\sqrt{-1}(1-x)F}\,Y\,e^{-\sqrt{-1}xF}dx\qquad(Y\in\mathcal C)
\]
are defined, with $\mathfrak d$ and $\mathfrak d^\star$ vanishing on the
exponentials and with the derivative of $\mathsf e$ given by the Duhamel
formula $\delta_\xi\mathsf e=-\sqrt{-1}\,\Xi_{\delta_\xi F}$. Then the one-form
$\vartheta(\xi)=\tau(\mathsf e\,\xi)$ on $\mathcal M$ is closed.
\end{lemma}

\begin{proof}
We use three identities, for $\mathfrak b\in\{\mathfrak d,\mathfrak d^\star\}$
and homogeneous $X,Y\in\mathcal C$:
\begin{enumerate}
\item[(I1)] $\tau\bigl((\mathfrak bX)\,Y\bigr)=-(-1)^{|X|}\tau\bigl(X\,\mathfrak bY\bigr)$,
 because $\mathfrak b(XY)=(\mathfrak bX)Y+(-1)^{|X|}X\,\mathfrak bY$ and
 $\tau\circ\mathfrak b=0$;
\item[(I2)] $\mathfrak b\,\Xi_Y=\Xi_{\mathfrak bY}$, because $\mathfrak b$
 vanishes on the exponentials of $F$;
\item[(I3)] $\tau(\Xi_YY')=\tau(Y\,\Xi_{Y'})$ for even $Y,Y'$, by the
 cyclicity of $\tau$ on even elements and the substitution $x\mapsto1-x$.
\end{enumerate}
Moreover $\frac d{dx}\bigl(e^{-\sqrt{-1}(1-x)F}Ye^{-\sqrt{-1}xF}\bigr)
=\sqrt{-1}\,[F,e^{-\sqrt{-1}(1-x)F}Ye^{-\sqrt{-1}xF}]$ gives
$[F,\Xi_Y]=\sqrt{-1}\,[\mathsf e,Y]$. From
$\delta_\xi F=-\mathfrak d\mathfrak d^\star\xi=[F,\xi]+\mathfrak d^\star\mathfrak d\xi$
and (I2),
\[
 \delta_\xi\mathsf e=-\sqrt{-1}\,[F,\Xi_\xi]-\sqrt{-1}\,\mathfrak d^\star\mathfrak d\,\Xi_\xi
 =[\mathsf e,\xi]-\sqrt{-1}\,\mathfrak d^\star\mathfrak d\,\Xi_\xi .
\]
For a two-parameter family,
\[
 \partial_v\vartheta(\xi)-\partial_u\vartheta(\eta)
 =\tau(\delta_\eta\mathsf e\;\xi)-\tau(\delta_\xi\mathsf e\;\eta)
 +\tau(\mathsf e\,[\xi,\eta]).
\]
The terms with the first summand of $\delta\mathsf e$, together with the last
term, give
\[
 \tau([\mathsf e,\eta]\xi)-\tau([\mathsf e,\xi]\eta)+\tau(\mathsf e[\xi,\eta])
 =-\tau(\mathsf e[\xi,\eta]),
\]
by cyclicity. For the remaining terms put
$\beta(\eta,\xi)=\tau(\Xi_\eta\,\mathfrak d\mathfrak d^\star\xi)$ and
$\beta'(\eta,\xi)=\tau(\Xi_\eta\,\mathfrak d^\star\mathfrak d\xi)$. By (I1),
applied twice,
\[
 \tau\bigl((\mathfrak d^\star\mathfrak d\,\Xi_\eta)\,\xi\bigr)
 =\tau\bigl((\mathfrak d\,\Xi_\eta)(\mathfrak d^\star\xi)\bigr)
 =-\tau\bigl(\Xi_\eta\,\mathfrak d\mathfrak d^\star\xi\bigr)=-\beta(\eta,\xi),
\]
and by (I2), (I3) and cyclicity,
\[
 \tau\bigl((\mathfrak d^\star\mathfrak d\,\Xi_\eta)\,\xi\bigr)
 =\tau\bigl(\Xi_{\mathfrak d^\star\mathfrak d\eta}\,\xi\bigr)
 =\tau\bigl((\mathfrak d^\star\mathfrak d\eta)\,\Xi_\xi\bigr)=\beta'(\xi,\eta).
\]
So $\beta(\eta,\xi)=-\beta'(\xi,\eta)$. Since
$\mathfrak d\mathfrak d^\star\xi+\mathfrak d^\star\mathfrak d\xi=-[F,\xi]$,
\[
 \beta(\eta,\xi)+\beta'(\eta,\xi)=-\tau(\Xi_\eta[F,\xi])=\tau([F,\Xi_\eta]\,\xi)
 =\sqrt{-1}\,\tau([\mathsf e,\eta]\xi)=-\sqrt{-1}\,\tau(\mathsf e[\xi,\eta]).
\]
Therefore
\begin{align*}
 -\sqrt{-1}\bigl[\tau\bigl((\mathfrak d^\star\mathfrak d\,\Xi_\eta)\xi\bigr)
 -\tau\bigl((\mathfrak d^\star\mathfrak d\,\Xi_\xi)\eta\bigr)\bigr]
 &=\sqrt{-1}\bigl[\beta(\eta,\xi)-\beta(\xi,\eta)\bigr]\\
 &=\sqrt{-1}\bigl[\beta(\eta,\xi)+\beta'(\eta,\xi)\bigr]=\tau(\mathsf e[\xi,\eta]),
\end{align*}
and the two contributions cancel.
\end{proof}

\begin{theorem}[Closedness]\label{thm:closed}
The one-form $\vartheta$ is closed. Hence there is a function $\cS_{\CC}$ on the
space of metrics, unique up to an additive constant, with
$d\cS_{\CC}(\xi)=\tr(\Omega\,\xi)$.
\end{theorem}

\begin{proof}
We apply Lemma~\ref{lem:abstract} to $\mathcal C=\mathcal B$, the functional
$\tau$ of Lemma~\ref{lem:tau}, the space of metrics, and
$\mathfrak d=[A,\,\cdot\,]$, $\mathfrak d^\star=[A^*,\,\cdot\,]$. The
functional vanishes on the images of these inner derivations by
Lemma~\ref{lem:tau}. Work in fixed bases. Adjoints depend on $h$ by
$T^\dagger=h^{-1}\overline T^{\mathsf T}h$, with $\overline T^{\mathsf T}$ the
conjugate transpose of the matrix, so for a variation $\xi$ one has
$\delta_\xi T^\dagger=[T^\dagger,\xi]$, hence $\delta_\xi A^*=[A^*,\xi]$ and
$\delta_\xi A=0$. The elements $A,A^*$ are odd, $\xi$ and $F$ are even. From
$A^2=0$, and hence $(A^*)^2=0$ for every metric, we get
$F_0:=(A-A^*)^2=-[A,A^*]$ and, by the graded Jacobi identity,
$[A,F_0]=0=[A^*,F_0]$ and
$\mathfrak d\mathfrak d^\star+\mathfrak d^\star\mathfrak d=[[A,A^*],\,\cdot\,]=-[F_0,\,\cdot\,]$.
The element $F_z=\sqrt{-1}\sum_if_i\otimes z_i\id_{E_i}$ is central in
$\mathcal B$ and does not depend on the metric. So $F=F_0+F_z$ satisfies
$\mathfrak dF=\mathfrak d^\star F=0$,
$\mathfrak d\mathfrak d^\star+\mathfrak d^\star\mathfrak d=-[F,\,\cdot\,]$ and
$\delta_\xi F=-[A,[A^*,\xi]]=-\mathfrak d\mathfrak d^\star\xi$. For a
two-parameter family of metrics with $\xi=h^{-1}\partial_uh$ and
$\eta=h^{-1}\partial_vh$ one has $\partial_v\xi-\partial_u\eta=[\xi,\eta]$.
The algebra is finite-dimensional, so the exponentials exist and the Duhamel
formula holds. Hence $\vartheta$ is closed, and the space of metrics is
contractible.
\end{proof}

\begin{theorem}[Primitive]\label{thm:primitive}
For $x\in[0,1]$ let $\Omega^{(x)}$ be the endomorphism \eqref{eq:Omegadef} of
the complex with differentials $\sqrt x\,d_i$ and the same arrows and metrics. Then
\[
 \cS_{\CC}(h)=\sum_n(-1)^n\Bigl[\sum_iz_i\log\det h_i^n-\sum_a\norm{\rho^n(a)}^2_{\HS}\Bigr]
 +\int_0^1\sum_{i,n}n\,\tr\bigl((\Omega^{(x)})_i^n-(\Omega^{(0)})_i^n\bigr)\frac{dx}x
\]
satisfies $d\cS_{\CC}(\xi)=\tr(\Omega\,\xi)$. Here $\rho^n(a)=\rho(a)|_{E_i^n}$,
$(\Omega^{(x)})_i^n$ is the block of $\Omega^{(x)}$ on $E_i^n$, and the
determinants are Gram determinants in fixed bases.
\end{theorem}

\begin{proof}
Rescaling the differentials by $\sqrt x$ at a fixed metric is the same, up
to unitary equivalence of the data, as the change of metric
$h_i^n\mapsto x^nh_i^n$ at fixed differentials: it multiplies each
$d_i^\dagger$ by $x$ and leaves $\rho(a)^\dagger$ unchanged. Along this path
of metrics the velocity is $\frac1x\Nop$, and $\Nop$ commutes with every
tangent vector $\xi$. By Theorem~\ref{thm:closed}, applied to the family
obtained from $he^{u\xi}$ by this rescaling,
$\delta_\xi\bigl[\frac1x\tr(\Omega^{(x)}\Nop)\bigr]=\partial_x\tr(\Omega^{(x)}\xi)$.
Since $\Omega^{(x)}$ is analytic in $x$, the integrand
$\frac1x\tr\bigl((\Omega^{(x)}-\Omega^{(0)})\Nop\bigr)$ is bounded and
$\delta_\xi$ of the integral is $\tr(\Omega^{(1)}\xi)-\tr(\Omega^{(0)}\xi)$;
here $\tr(\Omega^{(0)}\Nop)$ does not depend on the metric, because for zero
differentials $\Omega^{(0)}=\gr\,\Gamma$ and the arrow terms of
$\sum_i\tr\Gamma_i$ cancel in each degree. For zero differentials
$\Omega$ is that of a representation in each degree, up to the sign $\gr$,
with the primitive given by the first sum.
\end{proof}

For two-term complexes over $K_m$ with a surjective differential at the
vertex $0$ the formula of Theorem~\ref{thm:primitive} reduces to the
primitive of \cite[Proposition~7.1]{wue26a}. The integral in
Theorem~\ref{thm:primitive} is $\int_0^1\tr\bigl(\Nop\,(\Omega^{(x)}-\Omega^{(0)})\bigr)\frac{dx}x$:
a transgression by the number operator along the rescaling of the
differential, as in the construction of Bott--Chern forms of Hermitian
complexes by Bismut, Gillet and Soul\'e \cite{bgs88}, where the regularized
integral runs over $(0,\infty)$. Here the rescaling stops at $x=1$, and the
identity is on the space of metrics, in the algebra $\mathcal B$.

\section{The derivative of the phase potential along the flow}\label{sec:rate}

Assume $Z(\mathsf E)\neq0$ and fix $\varphi\in\RR$ with
$Z(\mathsf E)=|Z(\mathsf E)|e^{\ii\varphi}$. The \emph{phase potential} of the
metrized complex is $\Ima(e^{-\ii\varphi}\cS_{\CC})$, with $\cS_{\CC}$ as in
Theorem~\ref{thm:primitive}. Put
\[
 \cD:=\sum_{i,n}\cD_i^n,\qquad
 \cD_i^n:=\tr\Bigl(|\Omega_i^n|\,(\Phi_i^n-\varphi)\sin(\Phi_i^n-\varphi)\Bigr).
\]
We call $\cD_i^n$ the contribution of the block $\Omega_i^n$. By
Proposition~\ref{prop:rate} below, the derivative of the phase potential
along the flow is $-\cD$.

\begin{proposition}[Representations]\label{prop:kempfness}
Let $\mathsf E$ be a representation in degree zero, with $\varphi\in(0,\pi)$,
and put $\theta_i=\Ima(e^{-\ii\varphi}z_i)/\sin\varphi$ and
$\mu_i=\sum_{a:i\to\cdot}\rho(a)^\dagger\rho(a)-\sum_{a:\cdot\to i}\rho(a)\rho(a)^\dagger$,
the moment map. Then $\sum_i\theta_i\dim E_i=0$ and
\[
 \frac{\Ima(e^{-\ii\varphi}\cS_{\CC})}{\sin\varphi}
 =\sum_a\norm{\rho(a)}^2_{\HS}+\sum_i\theta_i\log\det h_i=:\cK_\theta(h),\qquad
 d\cK_\theta(\xi)=-\sum_i\tr\bigl(\xi_i(\mu_i-\theta_i)\bigr).
\]
This is the Kempf--Ness function of the representation for the character
$-\theta$ in the convention of King \cite{kin94}, for which a representation
is semistable when $\sum_i\theta_i\dim E'_i\leq0$ for all subrepresentations
$E'$, that is, when no subrepresentation has a central charge of argument
larger than $\varphi$. The metric is harmonic if and only if
$\mu_i=\theta_i$ for all $i$. All contributions $\cD_i^0$ are nonnegative,
and $\cD=0$ exactly at the harmonic metrics.
\end{proposition}

\begin{proof}
Here $\cS_{\CC}=\sum_iz_i\log\det h_i-\sum_a\norm{\rho(a)}^2_{\HS}$,
$\Ima(e^{-\ii\varphi}z_i)=\theta_i\sin\varphi$ and
$\Ima(e^{-\ii\varphi})=-\sin\varphi$; and
$\sum_i\theta_i\dim E_i=\Ima(e^{-\ii\varphi}Z(\mathsf E))/\sin\varphi=0$. The
block at $i$ is $z_i+\mu_i$, so
$\Ima\bigl(e^{-\ii\varphi}\tr(\Omega\xi)\bigr)=-\sin\varphi\sum_i\tr\bigl((\mu_i-\theta_i)\xi_i\bigr)$.
For a real number $x$, $\arg(z_i+x)=\varphi$ exactly for $x=\theta_i$. The
phases lie in $(0,\pi)$, so $\Phi_i-\varphi$ has its spectrum in
$(-\pi,\pi)$, where $x\sin x\geq0$ with equality only at $x=0$.
\end{proof}

\begin{proposition}\label{prop:rate}
Along the phase flow,
$\frac d{dt}\Ima\bigl(e^{-\ii\varphi}\cS_{\CC}\bigr)=-\cD$. The sum $\cD$ does not
depend on the choice of $\varphi$ modulo $2\pi$.
\end{proposition}

\begin{proof}
By Theorem~\ref{thm:closed} and the flow equation,
$\frac d{dt}\cS_{\CC}=-\tr(\Omega\Phi)$, and
$e^{-\ii\varphi}\tr\Omega=|Z(\mathsf E)|$ is real
(Proposition~\ref{prop:properties}(b)). Hence
$\frac d{dt}\Ima(e^{-\ii\varphi}\cS_{\CC})
=-\Ima\tr\bigl(e^{-\ii\varphi}\Omega(\Phi-\varphi)\bigr)$, and on each block
$e^{-\ii\varphi}\Omega_i^n(\Phi_i^n-\varphi)
=|\Omega_i^n|\,(\Phi_i^n-\varphi)e^{\ii(\Phi_i^n-\varphi)}$; the trace of a
product of two Hermitian matrices is real. Replacing $\varphi$ by
$\varphi+2\pi$ changes $\cD$ by
$-2\pi\sum\tr(|\Omega_i^n|\sin(\Phi_i^n-\varphi))=-2\pi\Ima(e^{-\ii\varphi}Z(\mathsf E))=0$.
\end{proof}

For a fixed $\varphi$ the summand $\cD_i^n$ is nonnegative when the spectrum
of $\Phi_i^n-\varphi$ lies in $\bigcup_{l\geq0}\pm[2l\pi,(2l+1)\pi]$. Since the
phase on $E_i^n$ lies in $(-\pi n,\pi(1-n))$, a block can have either sign as
soon as the complex has terms in more than one degree.

\subsection{A contractible summand that increases the potential}
The simplest instance of a negative contribution is a contractible summand
at one vertex.

\begin{proposition}\label{prop:contractible}
Let $\mathsf E=\mathsf E'\oplus\mathsf C$ be an orthogonal direct sum of
metrized complexes, where $\mathsf C=(C^n\xrightarrow{f}C^{n+1})$ is
concentrated at one vertex $i$ in degrees $n,n+1$, with $f$ invertible, and no
arrow acts on $\mathsf C$. The phase flow preserves the orthogonal splitting. The
blocks of $\Omega$ on $\mathsf C$ are $(-1)^nz_ie^{-f^\dagger f}$ and
$(-1)^{n+1}z_ie^{-ff^\dagger}$, with constant phases $\alpha_i-\pi n$ and
$\alpha_i-\pi(n+1)$, $\alpha_i=\arg z_i$; in orthonormal frames
$f(t)=e^{\pi t/2}f(0)$; and the contribution of $\mathsf C$ to $\cD$ is
\[
 \pi\,|z_i|\,\tr e^{-ff^\dagger}\,\sin(\alpha_i-\pi n-\varphi).
\]
\end{proposition}

\begin{proof}
For an orthogonal direct sum of metrized complexes $\Omega$ is the direct sum
of the endomorphisms of the summands, by Proposition~\ref{prop:coefficient}. On
$\mathsf C$, $\Delta=f^\dagger f\oplus ff^\dagger$, $\Gamma=z_i$ and all $M_a$
vanish. By \eqref{eq:frameflow},
$\dot f=\frac12(f\Phi^n-\Phi^{n+1}f)=\frac\pi2f$. With
$x=\alpha_i-\pi n-\varphi$ the two blocks contribute
$|z_i|\tr e^{-f^\dagger f}\,x\sin x+|z_i|\tr e^{-ff^\dagger}(x-\pi)\sin(x-\pi)
=\pi|z_i|\tr e^{-ff^\dagger}\sin x$.
\end{proof}

So the blocks of $\Omega$ on a contractible summand of this kind always tend
to zero, at the rate $e^{-\lambda e^{\pi t}}$ with $\lambda$ the least
eigenvalue of $f^\dagger f$ at time $0$; but the summand makes the potential
decrease only if $\sin(\alpha_i-\pi n-\varphi)>0$.

\begin{corollary}\label{cor:increase}
Let $Q=K_m$ with the charges $z_0=-1+\ii$, $z_1=1+\ii$. Let $\mathsf R$ be
a representation $(R_0\xrightarrow{\beta_j}R_1)$ with
$\dim R_0=\dim R_1\geq1$,
with a metric such that $\sum_j\beta_j^\dagger\beta_j=\id$ and
$\sum_j\beta_j\beta_j^\dagger=\id$, placed in degree $0$, and let
$\mathsf E=\mathsf R\oplus\mathsf C$ with $\mathsf C=(\CC\xrightarrow{f}\CC)$ at
the vertex $1$ in degrees $0,1$, $f\neq0$. Then $\varphi=\frac\pi2$, the
metric of $\mathsf R$ is harmonic and does not move modulo rescaling, and
along the whole trajectory
\[
 \frac d{dt}\Ima\bigl(e^{-\ii\varphi}\cS_{\CC}\bigr)=+\pi e^{-|f|^2}>0 .
\]
The phase potential is not bounded below on the metrics of $\mathsf E$.
\end{corollary}

\begin{proof}
The blocks of $\mathsf R$ are $z_0+\id$ and $z_1-\id$, both equal to
$\ii\,\id$, of phase $\frac\pi2=\varphi$. Proposition~\ref{prop:contractible}
with $i=1$, $n=0$, $\alpha_1=\frac\pi4$, $|z_1|=\sqrt2$ gives the rate
$-\pi\sqrt2\,e^{-|f|^2}\sin(\frac\pi4-\frac\pi2)$. By
Theorem~\ref{thm:primitive}, for $\mathsf C$ with metrics $h^0,h^1$ and
$b=|f|^2$ in orthonormal frames,
$\cS_{\CC}=z_1\bigl(\log h^0-\log h^1+\Ein(b)\bigr)$ with
$\Ein(b)=\int_0^1(1-e^{-xb})\,dx/x=\gamma+\log b+\Eone(b)$, where $\gamma$ is
Euler's constant and $\Eone(b)=\int_1^\infty e^{-xb}\,dx/x$ is the exponential
integral. Since $b$ is a constant multiple of $h^1/h^0$, this is
$z_1\Eone(b)$ up to a constant, and its contribution to
$\Ima(e^{-\ii\varphi}\cS_{\CC})$ is $-\Eone(b)$ up to a constant, which tends to
$-\infty$ as $b\to0$.
\end{proof}

For two-term complexes over $K_m$ with the acyclic part at the vertex $0$ in
degrees $0,1$ and cohomology that is nonzero at the vertex $1$, as in
\cite{wue26a,wue26b}, one has $\sin(\alpha_0-\varphi)>0$ for strictly ordered
charges. In general, for cohomology $\mathsf R[-n_*]$ with
$\mathsf R$ a representation of phase $\varphi_{\mathsf R}\in(0,\pi)$, so
that $\varphi=\varphi_{\mathsf R}-\pi n_*$, a contractible summand at the
vertex $i$ in degrees $n,n+1$ decreases the potential if and only if
$(-1)^{n-n_*}\sin(\alpha_i-\varphi_{\mathsf R})>0$.

\section{Duality}\label{sec:duality}

Let $\mathsf E$ be a metrized complex of representations of $Q$ with charges
$z_i$. The \emph{dual complex} $\mathsf E^\vee$ is the complex of
representations of the opposite quiver $Q^{\mathrm{op}}$ with
\[
 (E^\vee)_i^n=(E_i^{-n})^*,\qquad d_i^\vee=d_i^{\mathsf T},\qquad
 \rho^\vee(a^{\mathrm{op}})=\rho(a)^{\mathsf T},
\]
the dual metrics, and the charges $z_i^\vee=-\overline{z_i}$, which again lie
in the upper half-plane. Here $X^{\mathsf T}$ is the transposed map between
the dual spaces; in orthonormal frames and their dual frames it is the
transposed matrix, and $\overline X$ denotes the entrywise complex conjugate.

\begin{theorem}[Duality]\label{thm:duality}
For every metrized complex $\mathsf E$:
\begin{enumerate}
\item[(a)] As endomorphisms of $E^*$,
 $\Omega^\vee=-(\Omega^\dagger)^{\mathsf T}$, $|\Omega^\vee|=|\Omega|^{\mathsf T}$ and
 $\Phi^\vee=\pi-\Phi^{\mathsf T}$; in dual orthonormal frames,
 $\Omega^\vee=-\overline\Omega$ and $\Phi^\vee=\pi-\overline\Phi$. Hence
 $\mass(\mathsf E^\vee)=\mass(\mathsf E)$ and
 $Z(\mathsf E^\vee)=-\overline{Z(\mathsf E)}$.
\item[(b)] If $h(t)$ is a trajectory of the phase flow of $\mathsf E$, then
 $e^{-\pi t}(h(t)^{-1})^{\mathsf T}$ is a trajectory of the phase flow of
 $\mathsf E^\vee$. In orthonormal frames, transposition of all structure maps
 intertwines the equations \eqref{eq:frameflow} of the two complexes.
\item[(c)] If $Z(\mathsf E)\neq0$ and $\varphi^\vee=\pi-\varphi$, the phase
 potentials $\Ima(e^{-\ii\varphi^\vee}\cS^\vee_{\CC})$ and
 $\Ima(e^{-\ii\varphi}\cS_{\CC})$ agree, up to an additive constant, on
 corresponding metrics $h$ and $c\,(h^{-1})^{\mathsf T}$, $c>0$, and the
 contributions $\cD_i^n$ of corresponding blocks agree.
\end{enumerate}
\end{theorem}

\begin{proof}
(a) Use Proposition~\ref{prop:coefficient} for $\mathsf E^\vee$ in dual
orthonormal frames, where $(d^\vee)^\dagger=\overline d$ and
$(\rho^\vee)^\dagger=\overline\rho$. Then
$\Delta_i^\vee=d_i^{\mathsf T}\overline{d_i}+\overline{d_i}d_i^{\mathsf T}=\Delta_i^{\mathsf T}$. The arrows
of $Q^{\mathrm{op}}$ out of $i$ are the arrows of $Q$ into $i$, so
\[
 \Gamma_i^\vee=-\overline{z_i}+\Bigl(\sum_{a:k\to i}\rho(a)\rho(a)^\dagger
 -\sum_{a:i\to j}\rho(a)^\dagger\rho(a)\Bigr)^{\mathsf T}=-\overline{\Gamma_i},
\]
the bracket being Hermitian. For $a:i\to j$ the arrow $a^{\mathrm{op}}$ goes
from $j$ to $i$, and
$M^\vee_{a^{\mathrm{op}}}=\overline{d_i}\,\rho(a)^{\mathsf T}-\rho(a)^{\mathsf T}\overline{d_j}=-M_a^{\mathsf T}$,
$(M^\vee_{a^{\mathrm{op}}})^\dagger=-\overline{M_a}$. In
Proposition~\ref{prop:coefficient} for $Q^{\mathrm{op}}$ the term of
$a^{\mathrm{op}}$ at the vertex $i$ is of incoming type,
\[
 +\mathcal J_{\Delta_i^{\mathsf T},\Delta_j^{\mathsf T}}\bigl(M_a^{\mathsf T},\overline{M_a}\bigr)
 =\bigl(\mathcal J_{\Delta_i,\Delta_j}(M_a^\dagger,M_a)\bigr)^{\mathsf T}
 =\overline{\mathcal J_{\Delta_i,\Delta_j}(M_a^\dagger,M_a)},
\]
which enters $\mathsf E$ at $i$ with the sign $-$; likewise the term at $j$
is $-\overline{\mathcal J_{\Delta_j,\Delta_i}(M_a,M_a^\dagger)}$, which
enters $\mathsf E$ with $+$. Finally
$K_{\Delta^{\mathsf T}}(-\overline\Gamma)=-\bigl(K_\Delta(\Gamma^\dagger)\bigr)^{\mathsf T}
=-\overline{K_\Delta(\Gamma)}$. Altogether
$\gr\,\Omega^\vee=-\overline{\gr\,\Omega}$, and the degrees $n$
and $-n$ have the same parity. If $\Omega=|\Omega|e^{\ii\Phi}$ on a block of
degree $n$, then
$-\overline\Omega=\overline{|\Omega|}\,e^{\ii(\pi-\overline\Phi)}$ with
$\overline{|\Omega|}>0$, and the spectrum of $\pi-\overline\Phi$ lies in
$(\pi n,\pi(1+n))$, the branch of degree $-n$.

(b) Transposing \eqref{eq:frameflow} and using $\Phi^{\mathsf T}=\overline\Phi$ gives
$\frac d{dt}T^{\mathsf T}=\frac12(\overline{\Phi_{\mathrm{source}}}\,T^{\mathsf T}-T^{\mathsf T}\,\overline{\Phi_{\mathrm{target}}})$.
In $\mathsf E^\vee$ the map $T^{\mathsf T}$ goes from the dual of the target to the dual
of the source, with phases $\pi-\overline{\Phi_{\mathrm{target}}}$ and
$\pi-\overline{\Phi_{\mathrm{source}}}$, and its equation
\eqref{eq:frameflow} is the same expression. For metrics: if
$g=c\,(h^{-1})^{\mathsf T}$, then $g^{-1}\dot g=\dot c/c+\Phi^{\mathsf T}$, and
$-\Phi^\vee=\Phi^{\mathsf T}-\pi$ gives $c=e^{-\pi t}$ up to a constant.

(c) For corresponding variations $\xi^\vee=-\overline\xi$,
$\tr(\Omega^\vee\xi^\vee)=\overline{\tr(\Omega\xi)}$, so
$\cS^\vee_{\CC}=\overline{\cS_{\CC}}$ up to a constant on the metrics
$h\leftrightarrow(h^{-1})^{\mathsf T}$, and
$\Ima(e^{-\ii(\pi-\varphi)}\overline{\cS_{\CC}})=\Ima(e^{-\ii\varphi}\cS_{\CC})$.
For a common factor $c>0$ the endomorphism $\Omega^\vee$ is unchanged and
$\cS^\vee_{\CC}(cg)=\cS^\vee_{\CC}(g)+Z(\mathsf E^\vee)\log c$; since
$e^{-\ii\varphi^\vee}Z(\mathsf E^\vee)$ is real, the phase potential is
invariant under common rescaling.
\end{proof}

\begin{remark}[Injective differentials]\label{rem:injective}
Duality exchanges complexes that are $V\xrightarrow{D}U$ in degrees $0,1$ at
a vertex $i_0$, with $D$ surjective, and a vector space in degree zero at the
other vertices, with complexes that are $U'\xrightarrow{D'}V'$ in degrees
$-1,0$ at $i_0$, with $D'$ injective. The cohomology representation of the
dual is the dual of the cohomology representation, with the dual metrics
(the quotient metric on $V'/D'U'$ is the metric of the orthogonal complement
of $D'U'$); subrepresentations correspond to quotients and arguments $\psi$
of central charges to $\pi-\psi$, so semistability, stability and
polystability correspond. By Theorem~\ref{thm:duality}, every statement
about the flow of the first kind of complex therefore gives one for the
second kind, with the charges $-\overline{z_i}$, the argument $\pi-\varphi$
of the central charge, and the metrics $e^{-\pi t}(h^{-1})^{\mathsf T}$; in
particular the rescaled metrics $e^{\varphi t}h(t)$ of the complex
correspond to the rescaled metrics $e^{(\pi-\varphi)t}h^\vee(t)$ of the dual,
so that the flow converges on the cohomology for one exactly when it does
for the other, and the acyclic part decays for one exactly when it does for
the other. This is how Theorem~\ref{thm:onevertex}(c) follows from
Theorem~\ref{thm:onevertex}(b), and how the decay and convergence theorems
of \cite{wue26b,wue26c} pass to complexes with an injective differential.

For $Q=K_m$ one has $Q^{\mathrm{op}}\cong Q$ by exchanging the two vertices,
and the dual charges are $(-\overline{z_1},-\overline{z_0})$. The pair
$(-1+\ii,1+\ii)$ is self-dual, and strictly ordered charges are mapped to
strictly ordered charges. For these self-dual charges the sufficient
conditions of Theorem~\ref{thm:onevertex}(b) and (c) cover different
dimension vectors when the acyclic part has rank at least two: for
cohomology of dimension vector $(k,l)$, part (b) applies to a surjective
differential at the vertex $0$ when
$\arg((l-k)+\ii(k+l))\geq\frac{\pi}3$, and part (c) to an injective
differential at the vertex $1$ when this argument is at most
$\frac{2\pi}3$; for $k=l$ both apply.
\end{remark}

\section{Where the phase potential is monotone}\label{sec:monotone}

By Corollary~\ref{cor:increase} the phase potential is not monotone on all
complexes. This section treats two families: complexes with cohomology in
degree zero whose acyclic part is at one vertex of an arbitrary quiver,
where the monotonicity known for $K_m$ \cite[Theorem~7.6]{wue26b} persists,
and complexes over $K_m$ with acyclic parts at both vertices, where it
persists for ranks at most one and fails in general for larger ranks
(Theorem~\ref{thm:twosidedfailure}, Corollary~\ref{cor:classification}). We
use $\cD$ and $\cD_i^n$ from Section~\ref{sec:rate}, $K_\Delta$, $\mathcal J$
and the divided differences $\mathsf f[\,\cdot\,]$ from
Section~\ref{sec:setting}, $J_L(M)$, $M_\pm$ and the inequality
\eqref{eq:kernel} from Theorem~\ref{thm:imported}, and the terms of
Definition~\ref{def:phase}.

\subsection{An acyclic part at one vertex}
Let $Q$ be a finite quiver, $i_0$ a vertex without loops, and consider
complexes $\mathsf E$ with $E_i=E_i^0$ for $i\neq i_0$ and
$E_{i_0}=V\oplus U$ in degrees $0,1$, with $d_{i_0}=D:V\to U$ surjective and
$r=\dim U\geq1$. The chain-map condition says that every arrow into $i_0$
maps into $\ker D$, and is empty for the arrows out of $i_0$. The cohomology
is a representation $\mathcal H$ of $Q$ in degree $0$, with $\ker D$ at $i_0$;
we assume $\mathcal H\neq0$. Then
$Z=Z(\mathsf E)=\sum_iz_i\dim\mathcal H_i=|Z|e^{\ii\varphi}$ with
$\varphi\in(0,\pi)$. Let $\sigma=\sin\varphi$, $\alpha_i=\arg z_i$,
$c_0=|z_{i_0}|\sin(\alpha_{i_0}-\varphi)$, $B=D^\dagger D$, $L=DD^\dagger$,
\[
 C_{\mathrm{out}}=\sum_{a:i_0\to j}\rho(a)^\dagger\rho(a),\qquad
 G_{\mathrm{in}}=\sum_{a:k\to i_0}\rho(a)\rho(a)^\dagger,\qquad
 J=J_L(DC_{\mathrm{out}}D^\dagger),
\]
and $R(b)=(1-e^{-b})/b$.

\begin{theorem}[Acyclic part at one vertex]\label{thm:onevertex}
For the complexes just described:
\begin{enumerate}
\item[(a)] The blocks of $\Omega$ are
 \[
 \Omega_V=z_{i_0}e^{-B}+K_B(C_{\mathrm{out}})-G_{\mathrm{in}},\qquad
 \Omega_U=J-z_{i_0}e^{-L},
 \]
 and for $j\neq i_0$
 \[
 \Omega_j=z_j+\sum_{a:j\to\cdot}\rho(a)^\dagger\rho(a)
 -\sum_{\substack{a:k\to j\\k\neq i_0}}\rho(a)\rho(a)^\dagger
 -\sum_{a:i_0\to j}\rho(a)R(B)\rho(a)^\dagger .
 \]
\item[(b)] If $\alpha_{i_0}\geq\varphi$, then for every metric
 \[
 \cD\geq\varphi c_0\tr e^{-L}+\sigma\bigl(\varphi\tr J-\pi\tr J_-\bigr)
 \geq\varphi c_0\tr e^{-L}+\sigma\bigl(\varphi-\tfrac{2\pi}5\bigr)\tr J .
 \]
 Hence the phase potential is nonincreasing along the flow if
 $\alpha_{i_0}\geq\varphi$ and either $\varphi\geq\frac{2\pi}5$ or $r=1$, and
 strictly decreasing if moreover $\alpha_{i_0}>\varphi$.
\item[(c)] Let instead $E_i=E_i^0$ for $i\neq i_0$ and
 $E_{i_0}=U\oplus V$ in degrees $-1,0$, with $d_{i_0}=D:U\to V$ injective,
 $r=\dim U\geq1$, and nonzero cohomology, so that
 $\varphi=\arg Z(\mathsf E)\in(0,\pi)$. If $\alpha_{i_0}\leq\varphi$ and
 either $\varphi\leq\frac{3\pi}5$ or $r=1$, the phase potential is
 nonincreasing along the phase flow.
\end{enumerate}
\end{theorem}

\begin{proof}
(a) Proposition~\ref{prop:coefficient}. For $a:i_0\to j$,
$M_a=-\rho(a)D^\dagger:U\to E_j$. For $a:k\to i_0$, $M_a=d_{i_0}^\dagger\rho(a)=0$,
because $\rho(a)$ takes values in $V$, on which $d_{i_0}^\dagger$ vanishes.
For arrows not meeting $i_0$, $M_a=0$. So at $i_0$, $\Delta=B\oplus L$ and
$\Gamma=(z_{i_0}+C_{\mathrm{out}}-G_{\mathrm{in}})\oplus z_{i_0}$;
$K_B(G_{\mathrm{in}})=G_{\mathrm{in}}$ because $G_{\mathrm{in}}$ is supported on
$\ker B$; and $M_a^\dagger M_a=D\rho(a)^\dagger\rho(a)D^\dagger$ is supported on
$U$, with $\Delta_j=0$. At $j\neq i_0$, $\Delta_j=0$ and the mixed term is
$\rho(a)D^\dagger\bigl[\int_0^1(1-x)e^{-xL}dx\bigr]D\rho(a)^\dagger
=\rho(a)(1-R(B))\rho(a)^\dagger$.

(b) All blocks other than $U$ are in degree $0$, with phases in $(0,\pi)$,
and contribute nonnegative terms to $\cD$. Let $(u_k)$ be an orthonormal
eigenbasis of $\Phi_U$ with eigenvalues $\psi_k\in(-\pi,0)$,
$q_k=\inner{u_k}{|\Omega_U|u_k}$, $\iota_k=\inner{u_k}{Ju_k}$ and
$e_k=\inner{u_k}{e^{-L}u_k}>0$. Since
$\inner{u_k}{\Omega_Uu_k}=q_ke^{\ii\psi_k}$,
$q_k\sin(\psi_k-\varphi)=\Ima(e^{-\ii\varphi}\inner{u_k}{\Omega_Uu_k})
=-(\sigma\iota_k+c_0e_k)$. If $\sigma\iota_k+c_0e_k\geq0$, then
$\psi_k-\varphi\in[-\pi,-\varphi)$ and the summand
$q_k(\psi_k-\varphi)\sin(\psi_k-\varphi)$ of the contribution of $U$ is at
least $\varphi(\sigma\iota_k+c_0e_k)$; otherwise
$\psi_k-\varphi\in(-\pi-\varphi,-\pi)$ and the summand is at least
$(\pi+\varphi)(\sigma\iota_k+c_0e_k)$. In both cases it is at least
$\varphi(\sigma\iota_k+c_0e_k)-\pi\sigma(\iota_k)_-$, because $c_0e_k\geq0$;
here $(\iota_k)_-=\max(-\iota_k,0)$. Summing, with
$\sum_k(\iota_k)_-\leq\tr J_-$, gives the first inequality; the second is
\eqref{eq:kernel}, since $DC_{\mathrm{out}}D^\dagger\geq0$. For $r=1$,
$J\geq0$. If $\alpha_{i_0}>\varphi$, then $c_0>0$ and $\cD>0$. (c) follows from
(b) by Theorem~\ref{thm:duality}: the dual complex has an acyclic part at
$i_0$ in degrees $0,1$ with surjective differential, the charge
$-\overline{z_{i_0}}$ of argument $\pi-\alpha_{i_0}$ at $i_0$, and central
charge of argument $\pi-\varphi$ (Remark~\ref{rem:injective}).
\end{proof}

For $Q=K_m$, $i_0=0$ and cohomology at both vertices, $\varphi$ lies between
$\alpha_1$ and $\alpha_0$, and the case $\alpha_0>\alpha_1$ of (b) is
\cite[Proposition~7.4 and Theorem~7.6]{wue26b}; here equality and cohomology
at one vertex are
allowed. For $r\geq2$ a lower bound on $\varphi$ cannot be dispensed with:
for $K_2$, the charges $z_0=10^{-6}(-1+\ii)$, $z_1=1+\frac\ii{500}$ and
cohomology of dimension vector $(1,1)$, where $\varphi=0.0020\ldots$, there
is a complex with $r=2$ and a metric with $\cD<0$
\cite[Proposition~10.2]{wue26b}. The
hypothesis on $\alpha_{i_0}$ cannot be dropped: by
Proposition~\ref{prop:contractible} a contractible summand at $i_0$ in
degrees $0,1$ contributes
$\pi|z_{i_0}|\tr e^{-ff^\dagger}\sin(\alpha_{i_0}-\varphi)$ to $\cD$, and
Corollary~\ref{cor:increase} is an instance with a negative sign. Decay of
the acyclic part and convergence on the cohomology
(Definition~\ref{def:phase}(v)) for the complexes of
Theorem~\ref{thm:onevertex} are the subject of \cite{wue26b}, for $K_m$ with
a surjective differential at the vertex $0$, and of \cite{wue26c}, for
quivers without oriented cycles.

\subsection{Acyclic parts at both vertices of a generalized Kronecker quiver}
Let $Q=K_m$. Consider the three-term complexes
\begin{equation}\label{eq:twosided}
 (0\to U_1)\ \xrightarrow{\ D_1\ }\ (V\xrightarrow{A_j}W)\ \xrightarrow{\ D_0\ }\ (U_0\to0)
 \qquad\text{in degrees }-1,0,1,
\end{equation}
with $D_1:U_1\to W$ injective and $D_0:V\to U_0$ surjective; the chain-map
conditions are empty. The cohomology is the representation
$\ker D_0\to W/D_1U_1$ in degree $0$. Put $r_0=\dim U_0$, $r_1=\dim U_1$,
$n_0=\dim\ker D_0$, $n_1=\dim W-r_1$, assume $n_0+n_1\geq1$, and let
$Z=n_0z_0+n_1z_1=|Z|e^{\ii\varphi}$ with $\varphi\in(0,\pi)$,
$\sigma=\sin\varphi$, $c_0=|z_0|\sin(\alpha_0-\varphi)$,
$c_1=|z_1|\sin(\varphi-\alpha_1)$, and
\[
 B_0=D_0^\dagger D_0,\quad L_0=D_0D_0^\dagger,\quad P_1=D_1D_1^\dagger,\quad
 L_1=D_1^\dagger D_1,\quad C=\sum_jA_j^\dagger A_j,\quad G=\sum_jA_jA_j^\dagger .
\]
Here $D_0$ is the differential at the vertex $0$ and $D_1$ the differential
at the vertex $1$; the subscripts are vertices, not degrees. For $r_1=0$
these are complexes as in Theorem~\ref{thm:onevertex}(b) with $i_0=0$, and
for $r_0=0$ as in Theorem~\ref{thm:onevertex}(c) with $i_0=1$.

\begin{proposition}\label{prop:twosidedcoefficient}
The blocks of $\Omega$ for a complex \eqref{eq:twosided} are
\begin{align*}
 \Omega_V&=z_0e^{-B_0}+K_{B_0}(C)-\sum_j\mathcal J_{B_0,L_1}(A_j^\dagger D_1,D_1^\dagger A_j),\\
 \Omega_W&=z_1e^{-P_1}-K_{P_1}(G)+\sum_j\mathcal J_{P_1,L_0}(A_jD_0^\dagger,D_0A_j^\dagger),\\
 \Omega_{U_0}&=\mathcal J_{U_0}-z_0e^{-L_0},\qquad
 \mathcal J_{U_0}:=\sum_j\mathcal J_{L_0,P_1}(D_0A_j^\dagger,A_jD_0^\dagger),\\
 \Omega_{U_1}&=-z_1e^{-L_1}-\mathcal J_{U_1},\qquad
 \mathcal J_{U_1}:=\sum_j\mathcal J_{L_1,B_0}(D_1^\dagger A_j,A_j^\dagger D_1),
\end{align*}
with phases in $(0,\pi)$ on $V,W$, in $(-\pi,0)$ on $U_0$ and in $(\pi,2\pi)$
on $U_1$. If the charges are ordered, then for every metric
\[
 \cD\geq\varphi c_0\tr e^{-L_0}+\sigma\bigl(\varphi\tr\mathcal J_{U_0}-\pi\tr(\mathcal J_{U_0})_-\bigr)
 +(\pi-\varphi)c_1\tr e^{-L_1}+\sigma\bigl((\pi-\varphi)\tr\mathcal J_{U_1}-\pi\tr(\mathcal J_{U_1})_-\bigr).
\]
\end{proposition}

\begin{proof}
Proposition~\ref{prop:coefficient} with $E_0=V\oplus U_0$, $d_0=D_0$,
$E_1=U_1\oplus W$, $d_1=D_1$, $\Delta_0=B_0\oplus L_0$,
$\Delta_1=L_1\oplus P_1$, and the operator
$M_{a_j}=d_1^\dagger\rho(a_j)-\rho(a_j)d_0^\dagger$ of \eqref{eq:DeltaMGamma},
which maps $V\to U_1$ by $D_1^\dagger A_j$ and $U_0\to W$ by
$-A_jD_0^\dagger$. The bound for the block $U_0$ is the argument of
Theorem~\ref{thm:onevertex}(b) with $J$ replaced by the Hermitian matrix
$\mathcal J_{U_0}$. For $U_1$, with $(u_k)$ an eigenbasis of $\Phi_{U_1}$,
eigenvalues $\psi_k\in(\pi,2\pi)$, $q_k=\inner{u_k}{|\Omega_{U_1}|u_k}$,
$\iota_k=\inner{u_k}{\mathcal J_{U_1}u_k}$ and
$e_k=\inner{u_k}{e^{-L_1}u_k}$:
$q_k\sin(\psi_k-\varphi)=\sigma\iota_k+c_1e_k$. If this is nonnegative,
$\psi_k-\varphi\in(\pi-\varphi,\pi]$ and the summand is at least
$(\pi-\varphi)(\sigma\iota_k+c_1e_k)$; otherwise
$\psi_k-\varphi\in(\pi,2\pi-\varphi)$ and it is at least
$(2\pi-\varphi)(\sigma\iota_k+c_1e_k)$. In both cases it is at least
$(\pi-\varphi)(\sigma\iota_k+c_1e_k)-\pi\sigma(\iota_k)_-$.
\end{proof}

In eigenbases of $L_0$, with eigenvalues $\ell_i$, and of $P_1$, with
eigenvalues $p_\mu$ ($p_\mu=0$ on the orthogonal complement of $D_1U_1$),
\[
 (\mathcal J_{U_0})_{ik}=\sum_{j,\mu}(D_0A_j^\dagger)_{i\mu}\overline{(D_0A_j^\dagger)_{k\mu}}\;
 \mathsf f[\ell_i,p_\mu,\ell_k],
\]
with the second divided differences of $e^{-x}$. If all $p_\mu$ vanish this
is the Schur multiplier of \eqref{eq:kernel}; in general the interpolation
point $0$ is replaced by the eigenvalues $p_\mu$.

\begin{theorem}[Ranks at most one]\label{thm:rankone}
Let $\alpha_0\geq\alpha_1$ (ordered charges) and $r_0\leq1$, $r_1\leq1$. Then for every metric
on a complex \eqref{eq:twosided}
\[
 \cD\geq\varphi\bigl(c_0\tr e^{-L_0}+\sigma\tr\mathcal J_{U_0}\bigr)
 +(\pi-\varphi)\bigl(c_1\tr e^{-L_1}+\sigma\tr\mathcal J_{U_1}\bigr)\geq0,
\]
where a term with $r_i=0$ is omitted. The phase potential is nonincreasing
along the phase flow, and strictly decreasing if the charges are strictly
ordered, $n_0,n_1\geq1$ and $r_0+r_1\geq1$.
\end{theorem}

\begin{proof}
For $r_0=1$, $\mathcal J_{U_0}$ is the nonnegative number
\[
 \sum_j\int_{\substack{x,y\geq0\\x+y\leq1}}e^{-(x+y)L_0}\,D_0A_j^\dagger e^{-(1-x-y)P_1}A_jD_0^\dagger\,dx\,dy ;
\]
likewise $\mathcal J_{U_1}\geq0$ for $r_1=1$. Apply
Proposition~\ref{prop:twosidedcoefficient}. If $\alpha_0>\alpha_1$ and
$n_0,n_1\geq1$, then $c_0,c_1>0$.
\end{proof}

\begin{theorem}[Ranks $(2,1)$]\label{thm:twosidedfailure}
Let $z_0=-1+\ii$, $z_1=1+\ii$, and let $V=\CC^3$, $W=\CC^2$, $U_0=\CC^2$,
$U_1=\CC$ with the standard metrics,
\[
 D_0=\begin{pmatrix}0&1&0\\0&0&4\end{pmatrix},\qquad
 D_1=\begin{pmatrix}0\\3\end{pmatrix},\qquad
 A_1=\begin{pmatrix}1&0&0\\0&4&100\end{pmatrix},\qquad A_j=0\ (j\geq2).
\]
The cohomology is the representation $\CC\xrightarrow{1}\CC$, which is
stable, and the induced metric on it is harmonic. For these metrics
\[
 \cD=-1.7280352\ldots<0,\qquad\text{that is,}\qquad
 \frac d{dt}\Ima\bigl(e^{-\ii\varphi}\cS_{\CC}\bigr)=+1.728\ldots>0 .
\]
\end{theorem}

\begin{proof}
The complex is the orthogonal direct sum of the representation
$\CC\xrightarrow1\CC$, whose two blocks $z_0+1$ and $z_1-1$ equal $\ii$ and
have phase $\varphi=\frac\pi2$, and of the acyclic complex
$(0\to U_1)\to(\CC^2\xrightarrow{\rho}\CC)\to(U_0\to0)$ with
$\rho=(\rho_1,\rho_2)=(4,100)$ of Example~\ref{ex:nonnormal}. With
$(\ell_1,\ell_2)=(1,16)$ the eigenvalues of $B_0$ on $\CC^2$ and $p=9$,
Proposition~\ref{prop:twosidedcoefficient} gives for the acyclic complex the
blocks
\begin{align*}
 (\Omega_V)_{ik}&=z_0e^{-\ell_i}\delta_{ik}+\rho_i\rho_k\bigl(-\mathsf f[\ell_i,\ell_k]-p\,\mathsf f[\ell_i,p,\ell_k]\bigr),\\
 \Omega_W&=(z_1-|\rho|^2)e^{-p}+\sum_i\rho_i^2\ell_i\,\mathsf f[p,\ell_i,p],\\
 (-\Omega_{U_0})_{ik}&=z_0e^{-\ell_i}\delta_{ik}-\sqrt{\ell_i\ell_k}\,\rho_i\rho_k\,\mathsf f[\ell_i,p,\ell_k],\\
 -\Omega_{U_1}&=z_1e^{-p}+p\sum_i\rho_i^2\,\mathsf f[p,\ell_i,p],
\end{align*}
with the divided differences of $e^{-x}$. The four block contributions, on
$V$, $W$, $U_0$ and $U_1$, are $3.4635\ldots$, $2.0009\ldots$,
$-10.6257\ldots$ and $3.4332\ldots$; they and their sum were enclosed with
interval arithmetic \cite{joh17},\footnote{The program
\texttt{theorem\_6\_4\_dissipation.py}, which uses the library Arb, and its
output accompany this paper as supplementary files. The program also checks
that $\tr\Omega=0$ for the acyclic complex and compares the blocks, in
floating point, with an evaluation of the definition \eqref{eq:Omegadef},
with the exponential taken in $\mathcal A(K_1)\otimes\End(E)$.} with radius below
$10^{-18}$. The $2\times2$ polar decompositions are explicit, and the phases
on $U_0$ are $-3.1005\ldots$ and $-0.0347\ldots$
\end{proof}

\begin{corollary}\label{cor:classification}
For the charges $(-1+\ii,1+\ii)$ and $n_0=n_1\geq1$, the phase potential is
nonincreasing along the phase flow for all metrics on all complexes
\eqref{eq:twosided} with ranks $(r_0,r_1)$ if and only if $\min(r_0,r_1)=0$
or $r_0=r_1=1$.
\end{corollary}

\begin{proof}
Here $\varphi=\frac\pi2$ and $\alpha_1<\varphi<\alpha_0$. If $r_0=r_1=0$ the
complex is a representation and $\cD\geq0$ by
Proposition~\ref{prop:kempfness}. If
exactly one of $r_0,r_1$ vanishes, this is Theorem~\ref{thm:onevertex}(b) or
(c) for $Q=K_m$; if $r_0=r_1=1$ it is Theorem~\ref{thm:rankone}. Conversely, by Theorem~\ref{thm:duality}, which
exchanges $r_0$ and $r_1$ and fixes these charges, we may assume $r_0\geq2$
and $r_1\geq1$. Take the orthogonal direct sum of $n_0$ copies of
$\CC\xrightarrow1\CC$, of the acyclic complex of
Theorem~\ref{thm:twosidedfailure}, and of contractible summands at the vertex
$0$ in degrees $0,1$ and at the vertex $1$ in degrees $-1,0$, of ranks
$r_0-2$ and $r_1-1$, with differentials $x\,\id$. By
Proposition~\ref{prop:contractible} the latter contribute
$\pi(r_0+r_1-3)e^{-x^2}$ to $\cD$, which tends to $0$ as $x\to\infty$.
\end{proof}

\begin{remark}\label{rem:interpolation}
The example of Theorem~\ref{thm:twosidedfailure} comes from the following
computation. Let $r_0=2$, let $P_1$ have one nonzero eigenvalue $p$, and let
there be one nonzero arrow, which maps into the corresponding line. Then
$(\mathcal J_{U_0})_{ik}=v_i\overline{v_k}\,\mathsf f[\ell_i,p,\ell_k]$ for a
vector $v\in\CC^2$, and for a $2\times2$ matrix of this form the supremum over
$v$ of $\tr(\mathcal J_{U_0})_-/\tr\mathcal J_{U_0}$ is the positive part of
$\frac12\bigl(\mathsf f[\ell_1,p,\ell_2]\,(\mathsf f[\ell_1,p,\ell_1]\mathsf f[\ell_2,p,\ell_2])^{-1/2}-1\bigr)$,
attained for $|v_1|^2\mathsf f[\ell_1,p,\ell_1]=|v_2|^2\mathsf f[\ell_2,p,\ell_2]$.
As $\ell_1\to0$ and $\ell_2\to\infty$ this tends to
\[
 \frac12\Bigl(\frac{R(p)\,e^{p/2}}{\sqrt{(e^{-p}-1+p)/p^2}}-1\Bigr),
\]
which is $(\sqrt2-1)/2$ at $p=0$ and tends to infinity with $p$. So
\eqref{eq:kernel} has no analogue that is uniform in $p$.

\end{remark}

Whether the acyclic parts decay along the phase flow of the complexes
\eqref{eq:twosided}, in the sense of Definition~\ref{def:phase}(v), is open.

\section{Complexes of bundles on a curve}\label{sec:bundles}

The dictionary of \cite[Remark~5.14]{hkkp26} recalled in
Section~\ref{sec:setting} suggests the following formulation for two-term
complexes of holomorphic vector bundles on a curve, through the
superconnection of Quillen \cite{qui85}; the remark states the analogy
between the algebra of the quiver and the forms on a Riemann surface, and
the definition for bundles below is ours. A two-term complex $E^0\to E^1$ is
a
holomorphic triple in the sense of Bradlow and Garc\'ia-Prada \cite{bg96},
and for $E^0=\mathcal O$ a holomorphic pair \cite{bra91}; the equations
studied there are the vortex equations. No equation for the pair of metrics
is proposed here: the phases $\Phi^0$ and $\Phi^1$ defined below take values
in the disjoint intervals $(0,\pi)$ and $(-\pi,0)$, so that, as for
complexes of representations (Remark~\ref{rem:phase}), a metric of constant
phase exists only on a single bundle, where it is a Hermitian--Einstein
metric. The object of this section is the flow \eqref{eq:bundleflow}. We give
the endomorphism
$\Omega$, a primitive and the derivative of the phase potential along the
flow, show that the linearization of the flow is parabolic, and work out the
case of a section of a line bundle. In this section ${}^\dagger$ denotes the
adjoint for the metrics, also for endomorphism-valued forms; $K_\Delta$ and
the divided differences $\mathsf f[\,\cdot\,]$ of $e^{-x}$ (not the bundle
map $f$) are as in Section~\ref{sec:setting}, and $\Eone$ and $\Ein$ as in
Section~\ref{sec:rate}. The results of the earlier sections are cited where
they are used.

\subsection{The endomorphism \texorpdfstring{$\Omega$}{Omega}}
Let $X$ be a compact Riemann surface with a K\"ahler form $\omega$,
$\int_X\omega=1$, and $z$ a complex number with $\Ima z>0$. For a holomorphic
vector bundle $E$ put $Z(E)=z\rk E-\deg E$; for $z=\ii$ this is the usual
central charge $-\deg E+\ii\rk E$. A \emph{metrized two-term
complex} is $E^0\xrightarrow{f}E^1$ with $f$ a holomorphic bundle map and
Hermitian metrics $h^0,h^1$. Let $\nabla$ be the Chern connections, $F_{h^n}$
their curvatures,
\[
 \mathfrak k^n=\tfrac{\ii}{2\pi}\Lambda_\omega F_{h^n}\in\End(E^n),
 \qquad B=f^\dagger f,\quad L=ff^\dagger,
\]
and, in a local holomorphic coordinate $w$ with
$\omega=\frac{\ii}2g\,dw\wedge d\bar w$ and $f_w=\nabla_{\partial_w}f$,
\[
 \mathfrak q^0(T)=\frac{f_w^\dagger\,T\,f_w}{\pi g},\qquad
 \mathfrak q^1(T)=\frac{f_w\,T\,f_w^\dagger}{\pi g},
\]
which do not depend on the coordinate. The endomorphisms $\mathfrak k^n$ are
Hermitian, and $\int_X\tr\mathfrak k^n\,\omega=\deg E^n$. Let $\gr$ be the
grading of $E^0\oplus E^1$, let $\mathbb A=\nabla+\varkappa(f+f^\dagger)$ be
the superconnection on $E^0\oplus E^1$ with $\varkappa^2=-2\pi\ii$, and
\[
 \mathfrak F_h=\tfrac{\ii}{2\pi}\mathbb A^2-z\,\omega
 =(B\oplus L)+\tfrac{\ii\varkappa}{2\pi}\bigl(\nabla^{1,0}f+\nabla^{0,1}f^\dagger\bigr)
 +\tfrac{\ii}{2\pi}F_h-z\,\omega .
\]
We define $\Omega=\Omega^0\oplus\Omega^1\in\End(E^0)\oplus\End(E^1)$ by
\begin{equation}\label{eq:bundlecoefficient}
 \Omega\,\omega=\gr\,\bigl[\exp(-\mathfrak F_h)\bigr]^{(2)},
\end{equation}
the two-form part of the exponential in the algebra of
$\End(E^0\oplus E^1)$-valued forms with the sign rule of graded tensor
products. Up to the normalization $\varkappa$, $\mathbb A$ is Quillen's
superconnection of the Hermitian complex, and the supertrace of
$\exp(-\mathfrak F_h)$ is $e^{z\omega}$ times its Chern character form in the
normalization $\str\exp(-\frac{\ii}{2\pi}\mathbb A^2)$, whose part of degree
two represents $-c_1(E^0)+c_1(E^1)$. In \eqref{eq:bundlecoefficient} the
operator $\bar\partial+\varkappa f$, which does not depend on the metric,
takes the place of the element $A$ of Section~\ref{sec:setting}; so the
arrows correspond to forms of type $(0,1)$, whereas
\cite[Remark~5.14]{hkkp26} compares them with forms of type $(1,0)$.

\begin{proposition}\label{prop:bundleformula}
The blocks of $\Omega$ are
\[
 \Omega^0=K_B(z-\mathfrak k^0)+\mathcal J^0,\qquad
 \Omega^1=\mathcal J^1-K_L(z-\mathfrak k^1),
\]
with
$\mathcal J^0=\int e^{-xB}\,\mathfrak q^0\bigl(e^{-(1-x-y)L}\bigr)\,e^{-yB}\,dx\,dy$
and
$\mathcal J^1=\int e^{-xL}\,\mathfrak q^1\bigl(e^{-(1-x-y)B}\bigr)\,e^{-yL}\,dx\,dy$
over $\{x,y\geq0,\ x+y\leq1\}$. The endomorphisms $\mathcal J^n$ are
Hermitian, so $\Ima\Omega^0=(\Ima z)e^{-B}>0$ and
$\Ima(-\Omega^1)=(\Ima z)e^{-L}>0$ at every point of $X$. Moreover
\[
 \int_X\tr(\Omega)\,\omega=Z(E^0)-Z(E^1).
\]
\end{proposition}

\begin{proof}
The Duhamel expansion of $\exp(-(B\oplus L)-Y)$ with
$Y=\beta+\frac{\ii}{2\pi}F_h-z\omega$,
$\beta=\frac{\ii\varkappa}{2\pi}(\nabla^{1,0}f+\nabla^{0,1}f^\dagger)$, has
the two-form part
$K_{B\oplus L}(z\omega-\frac{\ii}{2\pi}F_h)+\int e^{-x(B\oplus L)}\beta e^{-(1-x-y)(B\oplus L)}\beta e^{-y(B\oplus L)}$,
and
$\beta\,T\,\beta=\frac{\ii}{2\pi}(dw\otimes f_w+d\bar w\otimes f_w^\dagger)\,T\,(dw\otimes f_w+d\bar w\otimes f_w^\dagger)$.
By the sign rule, on $E^0$:
$(d\bar w\otimes f_w^\dagger)\,T\,(dw\otimes f_w)=dw\wedge d\bar w\otimes f_w^\dagger Tf_w$,
and on $E^1$:
$(dw\otimes f_w)\,T\,(d\bar w\otimes f_w^\dagger)=-dw\wedge d\bar w\otimes f_wTf_w^\dagger$.
Since $\frac{\ii}{2\pi}dw\wedge d\bar w=\frac1{\pi g}\omega$, the two-form
parts are $[K_B(z-\mathfrak k^0)+\mathcal J^0]\omega$ on $E^0$ and
$[K_L(z-\mathfrak k^1)-\mathcal J^1]\omega$ on $E^1$; the latter is multiplied
by $-1$. Hermiticity is as in Proposition~\ref{prop:properties}(a). For the
integral, replace $f$ by $uf$, $u\in[0,1]$: the derivative in $u$ of the
supertrace of $[\exp(-\mathfrak F)]^{(2)}$ is exact, by the transgression
formula for superconnections \cite{qui85}, so the integral equals its value
at $u=0$, which is $z(\rk E^0-\rk E^1)-(\deg E^0-\deg E^1)$.
\end{proof}

The formula has the shape of $\Omega$ for a two-term complex of
representations, with $-\mathfrak k^0$ in the place of the arrow term. The
phases $\Phi^0$, with spectrum in $(0,\pi)$, and $\Phi^1$, with spectrum in
$(-\pi,0)$, are defined pointwise on $X$ as in
Proposition~\ref{prop:properties}(c), and the \emph{phase flow} is
\begin{equation}\label{eq:bundleflow}
 (h^n)^{-1}\partial_th^n=-\Phi^n(h),\qquad n=0,1 .
\end{equation}
For a single bundle, $E^1=0$, it is $h^{-1}\partial_th=-\arg(z-\mathfrak k_h)$,
with $z-\mathfrak k_h$ a normal endomorphism; its fixed points modulo
rescaling are the metrics with $\mathfrak k_h$ a scalar, the
Hermitian--Einstein metrics. Donaldson's heat flow \cite{don85} has velocity
proportional to $\mu-\mathfrak k_h$, with $\mu$ the slope; the present flow
has velocity $-\arg(z-\mathfrak k_h)$, and the two have the same stationary
points modulo rescaling. For a line bundle this is a flow of the metric
by an angle, of the same kind as the flow of Jacob and Yau for the deformed
Hermitian--Yang--Mills equation \cite[(1.2)]{jy17}, and the condition that
the phase be constant is of the type of the equation for $Z$-critical
connections \cite[Definition~2.22]{dms24}.

\subsection{The primitive and the derivative along the flow}
For a bundle $E$ with a reference metric let $M$ be the function on the
space of metrics with $dM(\xi)=\int_X\tr(\mathfrak k_h\xi)\,\omega$ and
$M(h^{\mathrm{ref}})=0$; this one-form is closed \cite{don85}, which is also
the case $f=0$, $E^1=0$ of the proposition below.

\begin{proposition}\label{prop:bundleprimitive}
On the space of smooth metrics the one-form
$\xi\mapsto\int_X\tr(\Omega\,\xi)\,\omega$, $\xi=h^{-1}\delta h$, is closed. A
primitive is
\begin{multline*}
 \cS_{\CC}(h)=z\int_X\Bigl(\log\det\bigl((h^{0,\mathrm{ref}})^{-1}h^0\bigr)
 -\log\det\bigl((h^{1,\mathrm{ref}})^{-1}h^1\bigr)\Bigr)\omega
 -M(h^0)+M(h^1)\\
 +\int_0^1\!\!\int_X\tr\bigl(\Omega_{[x]}^1+z-\mathfrak k^1\bigr)\,\omega\,\frac{dx}x,
\end{multline*}
where $\Omega_{[x]}^1$ is the block on $E^1$ of the endomorphism
\eqref{eq:bundlecoefficient} for $\sqrt x\,f$, the analogue of the
$\Omega^{(x)}$ of Theorem~\ref{thm:primitive}; for $x=0$ this block is
$-(z-\mathfrak k^1)$.
\end{proposition}

\begin{proof}
We apply Lemma~\ref{lem:abstract}. Let $\mathcal C$ be the algebra of smooth
$\End(E^0\oplus E^1)$-valued forms on $X$, graded by the form degree plus the
degree of the endomorphism, with the sign rule of graded tensor products, and
let $\tau(\alpha)=\int_X\str\alpha^{(2)}$, the integral of the supertrace of
the two-form part; it vanishes on supercommutators. Let
$\mathbb A''=\bar\partial+\varkappa f$ and
$\mathbb A'=\nabla^{1,0}+\varkappa f^\dagger$ be the two parts of the
superconnection, and put
\[
 \mathfrak d=(2\pi)^{-1/2}[\mathbb A'',\,\cdot\,],\qquad
 \mathfrak d^\star=-(2\pi)^{-1/2}[\mathbb A',\,\cdot\,],\qquad
 F=\tfrac1{2\pi}\mathbb A^2+\sqrt{-1}\,z\,\omega=-\sqrt{-1}\,\mathfrak F_h .
\]
These are odd derivations of $\mathcal C$: first-order operators acting on
endomorphism-valued forms by supercommutators. One has
$(\mathbb A'')^2=0$, because $\bar\partial^2=0$, $\bar\partial f=0$ and
$f^2=0$, and $(\mathbb A')^2=0$ likewise; so
$\mathbb A^2=[\mathbb A',\mathbb A'']$, and by the graded Jacobi identity
$\mathfrak d$ and $\mathfrak d^\star$ annihilate $\mathbb A^2$ and satisfy
$\mathfrak d\mathfrak d^\star+\mathfrak d^\star\mathfrak d
=-\frac1{2\pi}[\mathbb A^2,\,\cdot\,]=-[F,\,\cdot\,]$; the form $\omega$ is
closed and central. By Stokes' theorem and the vanishing of the supertrace
on supercommutators, $\tau\circ\mathfrak d=0=\tau\circ\mathfrak d^\star$. For
a variation $\xi=h^{-1}\delta h$ of the metrics, $\delta_\xi\mathbb A''=0$ and
$\delta_\xi\mathbb A'=[\mathbb A',\xi]$ for the Chern connection and the
adjoint $f^\dagger$, hence
$\delta_\xi F=\frac1{2\pi}[\mathbb A'',[\mathbb A',\xi]]=-\mathfrak d\mathfrak d^\star\xi$;
and $\partial_v\xi-\partial_u\eta=[\xi,\eta]$ as before. The exponential of
$-\sqrt{-1}F=-\mathfrak F_h$ is the finite Duhamel expansion around the
form-degree zero part $-(B\oplus L)$, so the remaining hypotheses hold, and
$\tau(e^{-\sqrt{-1}F}\xi)=\int_X\tr(\Omega\,\xi)\,\omega$ by
\eqref{eq:bundlecoefficient}. All terms in the proof of the lemma are
integrals over $X$ of polynomials in $\xi$, $\eta$ and their covariant
derivatives up to order two, with smooth coefficients. The
primitive is obtained as in Theorem~\ref{thm:primitive} from the path of
metrics $(h^0,xh^1)$, $x\in(0,1]$, along which the velocity is $\frac1x$
times the number operator and which corresponds to the rescaling
$\sqrt x\,f$; the integrand of the $x$-integral is bounded near $x=0$ because
$\Omega_{[x]}$ is analytic in $x$. For $f=0$ one has
$\Omega=(z-\mathfrak k_{h^0})\oplus-(z-\mathfrak k_{h^1})$, with the primitive
given by the first three terms.
\end{proof}

In the terms of the theory of Bott--Chern forms: the integral of the
form-valued one-form $\xi\mapsto\str\bigl(\xi\,e^{-\mathfrak F_h}\bigr)$ along
a path of metrics from $h^{\mathrm{ref}}$ to $h$ is the secondary form of the
pair $(h^{\mathrm{ref}},h)$ for the character form
$\str\exp(-\mathfrak F_h)$, and $\cS_{\CC}(h)-\cS_{\CC}(h^{\mathrm{ref}})$ is
the integral over $X$ of its part of degree two. As in Section~\ref{sec:rate}, for
$Z=Z(E^0)-Z(E^1)\neq0$ and $Z=|Z|e^{\ii\varphi}$ the function
$\Ima(e^{-\ii\varphi}\cS_{\CC})$ is called the \emph{phase potential}.

For a single bundle of slope $\mu=\deg E/\rk E$, with
$Z(E)=|Z|e^{\ii\varphi}$, one has $\Ima(e^{-\ii\varphi}z)=-\mu\sin\varphi$, and
\[
 \frac{\Ima(e^{-\ii\varphi}\cS_{\CC})}{\sin\varphi}
 =M(h)-\mu\int_X\log\det\bigl((h^{\mathrm{ref}})^{-1}h\bigr)\,\omega .
\]
This is Donaldson's functional \cite{don85}, in the normalization given by
$\mathfrak k_h$; its critical points are the Hermitian--Einstein metrics
$\mathfrak k_h=\mu$.

\begin{proposition}\label{prop:bundlerate}
Let $Z=Z(E^0)-Z(E^1)$, let $\varphi$ be a real number with
$\Ima(e^{-\ii\varphi}Z)=0$, and let $h(t)$ be a smooth solution of
\eqref{eq:bundleflow} on an interval. Then
\[
 \frac d{dt}\Ima\bigl(e^{-\ii\varphi}\cS_{\CC}(h(t))\bigr)
 =-\int_X\sum_{n=0,1}\tr\Bigl(|\Omega^n|(\Phi^n-\varphi)\sin(\Phi^n-\varphi)\Bigr)\omega .
\]
\begin{enumerate}
\item[(a)] For a single bundle, with $\varphi=\arg Z(E)$, the right side is
 nonpositive and vanishes only at Hermitian--Einstein metrics.
\item[(b)] If $\rk E^0=\rk E^1$ and $\deg E^1>\deg E^0$, then $Z>0$.
 With $\varphi=0$ the right side is negative, so $\Ima\cS_{\CC}$ is
 strictly decreasing along the flow.
\end{enumerate}
\end{proposition}

\begin{proof}
Apply Proposition~\ref{prop:bundleprimitive} to the flow equation and
use the polar decomposition in each fibre as in Proposition~\ref{prop:rate}.
Under the integral one may replace $\Phi^n$ by $\Phi^n-\varphi$ in the
factor multiplying the sine, since the integral of
$\Ima(e^{-\ii\varphi}\tr\Omega)$ is $\Ima(e^{-\ii\varphi}Z)=0$ by
Proposition~\ref{prop:bundleformula}. (a) $\Phi$ and $\varphi$ lie
in $(0,\pi)$, and $\Phi=\varphi$ exactly when
$\mathfrak k_h=\mu$. (b) $x\sin x>0$ on $(0,\pi)$ and on $(-\pi,0)$.
\end{proof}

For $\rk E^0>\rk E^1$ one has $\varphi\in(0,\pi)$, the term of degree zero is
nonnegative and the sign of the term of degree one is not controlled in
general, as for complexes of representations; for $\rk E^0<\rk E^1$ the roles
are exchanged.

\subsection{Parabolicity}
A system $\partial_t\xi=\mathcal P\,\Delta_\omega\xi+{}$terms of lower order,
with $\mathcal P$ a smooth field of endomorphisms of the fibre and
$\Delta_\omega=-d^*d\leq0$ the Laplace--Beltrami operator of $\omega$, is
parabolic in the sense of Petrovskii if all eigenvalues of $\mathcal P$ have
positive real part at every point.

\begin{proposition}\label{prop:parabolic}
The linearization of \eqref{eq:bundleflow} at a smooth metric is
$\partial_t\xi^n=\frac1{4\pi}\mathcal P^n\,\Delta_\omega\xi^n+{}$terms of lower
order, where $\mathcal P^n=-\mathcal V^n\circ K^n$, $K^0=K_B$, $K^1=K_L$,
and $\mathcal V^n$ is the derivative of the phase $\Phi^n$ with respect to
the Hermitian part of $(-1)^n\Omega^n$ at fixed imaginary part.
All eigenvalues of $\mathcal P^n$ have positive real part: the system is
parabolic in the sense of Petrovskii. Moreover $\mathcal P^n$ has a positive definite
symmetric part for the fibre metric $\tr(K^n(\xi)\eta)$ on Hermitian
endomorphisms.
\end{proposition}

\begin{proof}
$\Omega$ depends on second derivatives of $h$ only through $\mathfrak k$, and
for a variation $\xi$ of the metric the variation of $\mathfrak k^n$ is
$-\frac1{4\pi}\Delta_\omega\xi^n$ up to terms of lower order. So the
second-order part of the variation of $(-1)^n\Omega^n$ is the Hermitian
endomorphism $\frac1{4\pi}K^n(\Delta_\omega\xi^n)$, with the imaginary part
unchanged, and the variation of the phase is $\mathcal V^n$ applied to it.
The operators $K_B$ and $K_L$ are positive definite and self-adjoint on
Hermitian matrices: in an eigenbasis of $B$, or of $L$, with eigenvalues
$\ell_i$, they multiply the entries by the positive numbers
$-\mathsf f[\ell_i,\ell_k]$. By
Theorem~\ref{thm:imported}(b), $\tr(\eta\,\mathcal V^n\eta)<0$ for Hermitian
$\eta\neq0$. Hence
$\tr\bigl(K(\xi)\,\mathcal P\xi\bigr)=-\tr\bigl(K(\xi)\,\mathcal V(K(\xi))\bigr)>0$,
which is the statement on the symmetric part for the inner product
$\tr(K(\xi)\eta)$; an operator with positive definite symmetric part for
some inner product has its eigenvalues in the open right half-plane.
\end{proof}

We do not prove short-time existence of smooth solutions of
\eqref{eq:bundleflow}. For a line bundle the flow
is the scalar equation
$\partial_tu=-\arg\bigl(z-\mathfrak k_{h^{\mathrm{ref}}}+\frac1{4\pi}\Delta_\omega u\bigr)$
for $h=e^uh^{\mathrm{ref}}$, whose right side is increasing in
$\Delta_\omega u$, and $\mathcal P=\Ima z/|z-\mathfrak k_h|^2$. This tends to
zero where the phase approaches $0$ or $\pi$, so the parabolicity is not
uniform.

\subsection{A section of a line bundle}
Let the complex be $\mathcal O\xrightarrow{s}\mathcal L$, with $\mathcal L$ a
line bundle of degree $d\geq1$ and $s$ a nonzero holomorphic section with
divisor $\Sigma$. It is quasi-isomorphic to $\mathcal O_\Sigma[-1]$, and
$Z=Z(\mathcal O)-Z(\mathcal L)=d>0$, so $\varphi=0$. Let
$\mathsf h=h^1/h^0$, a metric on $\mathcal L$, $b=|s|^2_{\mathsf h}$,
$\mathfrak k^n=\frac{\ii}{2\pi}\Lambda_\omega F_{h^n}$ as before and
$\nu=|\nabla_ws|^2_{\mathsf h}/(\pi g)\geq0$. All blocks are scalar, and
Proposition~\ref{prop:bundleformula} gives
\[
 \Omega^0=e^{-b}\bigl(z-\mathfrak k^0+\tfrac\nu2\bigr),\qquad
 \Omega^1=-e^{-b}\bigl(z-\mathfrak k^1-\tfrac\nu2\bigr),
\]
\[
 \Phi^0=\arg\bigl(z-\mathfrak k^0+\tfrac\nu2\bigr)\in(0,\pi),\qquad
 \Phi^1=\arg\bigl(z-\mathfrak k^1-\tfrac\nu2\bigr)-\pi\in(-\pi,0).
\]
The sum is
\[
 (\Omega^0+\Omega^1)\,\omega=e^{-|s|^2_{\mathsf h}}\bigl[\tfrac{\ii}{2\pi}F_{\mathsf h}
 +\tfrac{\ii}{2\pi}\nabla s\wedge\overline{\nabla s}\bigr]
 =\tfrac{\ii}{2\pi}F_{\mathsf h}+\tfrac{\ii}{2\pi}\partial\bar\partial\Ein\bigl(|s|^2_{\mathsf h}\bigr),
\]
a representative of the Euler class of $\mathcal L$ of the kind constructed
by Mathai and Quillen \cite{mq86}, pulled back by $s$; its integral is $d$,
and when $s$ is replaced by $cs$ with $c\to\infty$, it converges to the
current of integration over $\Sigma$, because
$\Ein(c^2b)=\gamma+\log c^2+\log b+\Eone(c^2b)$, with $\gamma$ Euler's
constant, and $\frac{\ii}{2\pi}\partial\bar\partial\log b=\delta_\Sigma-\frac{\ii}{2\pi}F_{\mathsf h}$
by the Poincar\'e--Lelong formula.

With $u^n=\log(h^n/h^{n,\mathrm{ref}})$ the flow is $\partial_tu^n=-\Phi^n$,
so $\partial_t\log b=\Phi^0-\Phi^1\in(0,2\pi)$: $b$ increases at every point
outside $\Sigma$. By Proposition~\ref{prop:bundleprimitive} the phase
potential is
\[
 \Ima\cS_{\CC}=(\Ima z)\,\mathcal W(h),\qquad
 \mathcal W(h)=\int_X\bigl(u^0-u^1+\Ein(b)\bigr)\omega
 =\int_X\Eone(b)\,\omega+\mathrm{const}.
\]
It involves no derivatives of the metrics, and it is bounded below, with a
bound that is not attained.

\begin{proposition}\label{prop:section}
Along every solution of the flow, at every point of $X$,
\[
 \partial_t\bigl(u^0-u^1+\Ein(b)\bigr)=-e^{-b}(\Phi^0-\Phi^1)<0,
 \qquad\frac d{dt}\mathcal W=-\int_Xe^{-b}\bigl(\Phi^0+|\Phi^1|\bigr)\,\omega .
\]
Consequently
$\int_0^T\!\int_Xe^{-b}(\Phi^0+|\Phi^1|)\,\omega\,dt\leq\int_X\Eone(b|_{t=0})\,\omega<\infty$
on the interval of existence.
\end{proposition}

\begin{proof}
$\partial_t(u^0-u^1)=-(\Phi^0-\Phi^1)$ and
$\partial_t\Ein(b)=\frac{1-e^{-b}}b\,\partial_tb=(1-e^{-b})(\Phi^0-\Phi^1)$.
The function $\Eone(b)$ has a logarithmic singularity along $\Sigma$ and is
integrable; $\mathcal W-\int_X\Eone(b)\,\omega$ is constant along the flow and
$\Eone\geq0$.
\end{proof}

So in this case the potential is monotone for every $z$ in the upper
half-plane, in accordance with Proposition~\ref{prop:bundlerate}(b). If
instead $d=0$ and $s$ has no zeros, the complex is acyclic and $Z=0$; for
constant metrics on a trivialization $\Omega^0=ze^{-b}=-\Omega^1$ and
$b(t)=b(0)e^{\pi t}$, as for the contractible summand of
Proposition~\ref{prop:contractible}.

\begin{thebibliography}{HKKP26}

\bibitem[BC65]{bc65}
R.~Bott and S.~S.~Chern,
\textit{Hermitian vector bundles and the equidistribution of the zeroes of
their holomorphic sections},
Acta Math. \textbf{114} (1965), 71--112.
\href{https://doi.org/10.1007/BF02391818}{doi:10.1007/BF02391818}.

\bibitem[BGS88]{bgs88}
J.-M.~Bismut, H.~Gillet and C.~Soul\'e,
\textit{Analytic torsion and holomorphic determinant bundles. I. Bott--Chern
forms and analytic torsion},
Comm. Math. Phys. \textbf{115} (1988), no.~1, 49--78.

\bibitem[Bra91]{bra91}
S.~B.~Bradlow,
\textit{Special metrics and stability for holomorphic bundles with global
sections},
J. Differential Geom. \textbf{33} (1991), no.~1, 169--213.

\bibitem[BG96]{bg96}
S.~B.~Bradlow and O.~Garc\'ia-Prada,
\textit{Stable triples, equivariant bundles and dimensional reduction},
Math. Ann. \textbf{304} (1996), no.~1, 225--252.
\href{https://arxiv.org/abs/alg-geom/9401008}{arXiv:alg-geom/9401008}.

\bibitem[Bri07]{bri07}
T.~Bridgeland,
\textit{Stability conditions on triangulated categories},
Ann. of Math. (2) \textbf{166} (2007), no.~2, 317--345.

\bibitem[DMS24]{dms24}
R.~Dervan, J.~B.~McCarthy and L.~M.~Sektnan,
\textit{$Z$-critical connections and Bridgeland stability conditions},
Camb. J. Math. \textbf{12} (2024), no.~2, 253--355.
\href{https://arxiv.org/abs/2012.10426v5}{arXiv:2012.10426v5}.

\bibitem[Don85]{don85}
S.~K.~Donaldson,
\textit{Anti self-dual Yang--Mills connections over complex algebraic
surfaces and stable vector bundles},
Proc. London Math. Soc. (3) \textbf{50} (1985), no.~1, 1--26.

\bibitem[HKKP26]{hkkp26}
F.~Haiden, L.~Katzarkov, M.~Kontsevich and P.~Pandit,
\textit{Towards categorical K\"ahler geometry},
preprint (2026), \href{https://arxiv.org/abs/2609.00978v1}{arXiv:2609.00978v1}.

\bibitem[JY17]{jy17}
A.~Jacob and S.-T.~Yau,
\textit{A special Lagrangian type equation for holomorphic line bundles},
Math. Ann. \textbf{369} (2017), no.~1--2, 869--898.
\href{https://arxiv.org/abs/1411.7457}{arXiv:1411.7457}.

\bibitem[Joh17]{joh17}
F.~Johansson,
\textit{Arb: efficient arbitrary-precision midpoint-radius interval
arithmetic},
IEEE Trans. Comput. \textbf{66} (2017), no.~8, 1281--1292.
\href{https://doi.org/10.1109/TC.2017.2690633}{doi:10.1109/TC.2017.2690633}.

\bibitem[KN79]{kn79}
G.~Kempf and L.~Ness,
\textit{The length of vectors in representation spaces},
in \textit{Algebraic geometry (Copenhagen, 1978)},
Lecture Notes in Math. \textbf{732}, Springer, Berlin, 1979, 233--243.

\bibitem[Kin94]{kin94}
A.~D.~King,
\textit{Moduli of representations of finite dimensional algebras},
Quart. J. Math. Oxford Ser. (2) \textbf{45} (1994), no.~4, 515--530.
\href{https://doi.org/10.1093/qmath/45.4.515}{doi:10.1093/qmath/45.4.515}.

\bibitem[MQ86]{mq86}
V.~Mathai and D.~Quillen,
\textit{Superconnections, Thom classes, and equivariant differential forms},
Topology \textbf{25} (1986), no.~1, 85--110.

\bibitem[Qui85]{qui85}
D.~Quillen,
\textit{Superconnections and the Chern character},
Topology \textbf{24} (1985), no.~1, 89--95.

\bibitem[Wue26a]{wue26a}
B.~J.~Wuebben,
\textit{Mass increase and acyclic cancellation for phase flows on quiver
complexes},
preprint (2026),
\href{https://doi.org/10.5281/zenodo.23921919}{doi:10.5281/zenodo.23921919}.

\bibitem[Wue26b]{wue26b}
B.~J.~Wuebben,
\textit{Convergence of the minimizing flow on Kronecker quiver complexes},
preprint (2026), \href{https://doi.org/10.5281/zenodo.24341605}{doi:10.5281/zenodo.24341605}.

\bibitem[Wue26c]{wue26c}
B.~J.~Wuebben,
\textit{Convergence of the minimizing flow on acyclic quivers},
preprint (2026), \href{https://doi.org/10.5281/zenodo.24346923}{doi:10.5281/zenodo.24346923}.

\end{thebibliography}

\end{document}
